% /*@@
%   @file      UtilReference.tex
%   @date      Sat Nov  3 18:47:53 MET 2001
%   @author    Jonathan Thornburg <jthorn@aei.mpg.de>
%   @desc
%   Utility Functions Reference for the Cactus User's Guide
%   (this was briefly part of FunctionReference.tex in this directory,
%   before that this was part of ../UsersGuide/FunctionReference.tex;
%   see there for older cvs history)
%   @enddesc
%   @version $Header$
% @@*/

\begin{cactuspart}{{\tt Util\_*} Functions Reference}{$RCSfile$}{$Revision$}
\label{part:UtilReference}
\renewcommand{\thepage}{\Alph{part}\arabic{page}}

%%%%%%%%%%%%%%%%%%%%%%%%%%%%%%%%%%%%%%%%%%%%%%%%%%%%%%%%%%%%%%%%%%%%%%%%%%%%%%%%
%%%%%%%%%%%%%%%%%%%%%%%%%%%%%%%%%%%%%%%%%%%%%%%%%%%%%%%%%%%%%%%%%%%%%%%%%%%%%%%%
%%%%%%%%%%%%%%%%%%%%%%%%%%%%%%%%%%%%%%%%%%%%%%%%%%%%%%%%%%%%%%%%%%%%%%%%%%%%%%%%

In this chapter all \hbox{{\tt Util\_}*} Cactus utility functions are
described.  These are low-level functions mainly for more complicated
programming, which are used by the rest of Cactus, but don't depend
heavily on it.  Some of them are callable from Fortran or C, but
many are C-only.

%%%%%%%%%%%%%%%%%%%%%%%%%%%%%%%%%%%%%%%%%%%%%%%%%%%%%%%%%%%%%%%%%%%%%%%%%%%%%%%%

\chapter{Functions Alphabetically}

Here the functions are listed alphabetically within each section.

%%%%%%%%%%%%%%%%%%%%%%%%%%%%%%%%%%%%%%%%

%% comment this whole section out until it's non-empty
%%\section{Miscellaneous Functions}
%%
%%\begin{Lentry}

%%\item[Util\_asprintf]
%%     [\pageref{Util-asprintf}]

%%\item[Util\_BinTreeFindNode]
%%     [\pageref{Util-BinTreeFindNode}]

%%\item[Util\_BinTreePrintNodes]
%%     [\pageref{Util-BinTreePrintNodes}]

%%\item[Util\_BinTreeStoreData]
%%     [\pageref{Util-BinTreeStoreData}]

%%\item[Util\_BinTreeTraverseInorder]
%%     [\pageref{Util-BinTreeTraverseInorder}]

%%\item[Util\_BinTreeTraversePostorder]
%%     [\pageref{Util-BinTreeTraversePostorder}]

%%\item[Util\_BinTreeTraversePreorder]
%%     [\pageref{Util-BinTreeTraversePreorder}]

%%\item[Util\_CacheMalloc]
%%     [\pageref{Util-CacheMalloc}]

%%\item[Util\_CurrentDate]
%%     [\pageref{Util-CurrentDate}]

%%\item[Util\_CurrentTime]
%%     [\pageref{Util-CurrentTime}]

%%\item[Util\_DeleteHandle]
%%     [\pageref{Util-DeleteHandle}]

%%\item[Util\_DoubleInRange]
%%     [\pageref{Util-DoubleInRange}]

%%\item[Util\_DoubleInRangeList]
%%     [\pageref{Util-DoubleInRangeList}]

%%\item[Util\_GetHandle]
%%     [\pageref{Util-GetHandle}]

%%\item[Util\_GetHandleName]
%%     [\pageref{Util-GetHandleName}]

%%\item[Util\_GetHandledData]
%%     [\pageref{Util-GetHandledData}]

%%\item[Util\_GetHostName]
%%     [\pageref{Util-GetHostName}]

%%\item[Util\_HashAdd]
%%     [\pageref{Util-HashAdd}]

%%\item[Util\_HashCreate]
%%     [\pageref{Util-HashCreate}]

%%\item[Util\_HashData]
%%     [\pageref{Util-HashData}]

%%\item[Util\_HashDelete]
%%     [\pageref{Util-HashDelete}]

%%\item[Util\_HashDestroy]
%%     [\pageref{Util-HashDestroy}]

%%\item[Util\_HashHash]
%%     [\pageref{Util-HashHash}]

%%\item[Util\_HashStore]
%%     [\pageref{Util-HashStore}]

%%\item[Util\_InList]
%%     [\pageref{Util-InList}]

%%\item[Util\_IntInRange]
%%     [\pageref{Util-IntInRange}]

%%\item[Util\_IntInRangeList]
%%     [\pageref{Util-IntInRangeList}]

%%\item[Util\_NewHandle]
%%     [\pageref{Util-NewHandle}]

%%\item[Util\_NullTerminateString]
%%     [\pageref{Util-NullTerminateString}]

%%\item[Util\_SplitFilename]
%%     [\pageref{Util-SplitFilename}]

%%\item[Util\_SplitString]
%%     [\pageref{Util-SplitString}]

%%\item[Util\_StringListAdd]
%%     [\pageref{Util-StringListAdd}]

%%\item[Util\_StringListCreate]
%%     [\pageref{Util-StringListCreate}]

%%\item[Util\_StringListDestroy]
%%     [\pageref{Util-StringListDestroy}]

%%\item[Util\_StringListNext]
%%     [\pageref{Util-StringListNext}]

%%\end{Lentry}

%%%%%%%%%%%%%%%%%%%%%%%%%%%%%%%%%%%%%%%%

\section{String Functions}

\begin{Lentry}

\item[Util\_StrCmpi]
     [\pageref{Util-StrCmpi}]
Compare two strings, ignoring upper/lower case.

\item[Util\_Strdup]
     [\pageref{Util-Strdup}]
``Duplicate'' a string,
\ie{} copy it to a newly--\verb|malloc|ed buffer.

\item[Util\_Strlcat]
     [\pageref{Util-Strlcat}]
Concatenate two strings safely.

\item[Util\_Strlcpy]
     [\pageref{Util-Strlcpy}]
Copy a string safely.

\end{Lentry}

%%%%%%%%%%%%%%%%%%%%%%%%%%%%%%%%%%%%%%%%

\section{Table Functions}

\begin{Lentry}

\item[Util\_TableClone]
     [\pageref{Util-TableClone}]
Create a new table which is a ``clone'' (exact copy) of an existing
table

\item[Util\_TableCreate]
     [\pageref{Util-TableCreate}]
Create a new (empty) table

\item[Util\_TableCreateFromString]
     [\pageref{Util-TableCreateFromString}]
Create a new table (with the case-insensitive flag set) and sets
values in it based on a string argument (interpreted with
``parameter-file'' semantics)

\item[Util\_TableDeleteKey]
     [\pageref{Util-TableDeleteKey}]
Delete a specified key/value entry from a table

\item[Util\_TableDestroy]
     [\pageref{Util-TableDestroy}]
Destroy a table

\item[Util\_TableGet*]
     [\pageref{Util-TableGet*}]
This is a family of functions, one for each Cactus data type,
to get the single (1-element array) value, or more generally the
first array element of the value, associated with a specified key
in a key/value table.

\item[Util\_TableGet*Array]
     [\pageref{Util-TableGet*Array}]
This is a family of functions, one for each Cactus data type,
to get a copy of the value associated with a specified key, and store
it (more accurately, as much of it as will fit) in a specified array

\item[Util\_TableGetGeneric]
     [\pageref{Util-TableGetGeneric}]
Get the single (1-element array) value, or more generally the
first array element of the value, associated with a specified key
in a key/value table; the value's data type is generic

\item[Util\_TableGetGenericArray]
     [\pageref{Util-TableGetGenericArray}]
Get a copy of the value associated with a specified key, and store
it (more accurately, as much of it as will fit) in a specified array;
the array's data type is generic

\item[Util\_TableGetString]
     [\pageref{Util-TableGetString}]
Gets a copy of the character-string value associated with a specified
key in a table, and stores it (more accurately, as much of it as will fit)
in a specified character string

\item[Util\_TableItAdvance]
     [\pageref{Util-TableItAdvance}]
Advance a table iterator to the next entry in the table

\item[Util\_TableItClone]
     [\pageref{Util-TableItClone}]
Creates a new table iterator which is a ``clone'' (exact copy) of an
existing table iterator

\item[Util\_TableItCreate]
     [\pageref{Util-TableItCreate}]
Create a new table iterator

\item[Util\_TableItDestroy]
     [\pageref{Util-TableItDestroy}]
Destroy a table iterator

\item[Util\_TableItQueryIsNonNull]
     [\pageref{Util-TableItQueryIsNonNull}]
Query whether a table iterator is {\em not\/} in the ``null-pointer'' state

\item[Util\_TableItQueryIsNull]
     [\pageref{Util-TableItQueryIsNull}]
Query whether a table iterator is in the ``null-pointer'' state

\item[Util\_TableItQueryKeyValueInfo]
     [\pageref{Util-TableItQueryKeyValueInfo}]
Query the key and the type and number of elements of the value
corresponding to that key, of the table entry to which an iterator points

\item[Util\_TableItQueryTableHandle]
     [\pageref{Util-TableItQueryTableHandle}]
Query what table a table iterator iterates over

\item[Util\_TableItResetToStart]
     [\pageref{Util-TableItResetToStart}]
Reset a table iterator to point to the starting table entry

\item[Util\_TableItSetToKey]
     [\pageref{Util-TableItSetToKey}]
Set a key/value iterator to point to a specified entry in the table.

\item[Util\_TableItSetToNull]
     [\pageref{Util-TableItSetToNull}]
Set a key/value iterator to the ``null-pointer'' state.

\item[Util\_TableQueryFlags]
     [\pageref{Util-TableQueryFlags}]
Query a table's flags word

\item[Util\_TableQueryValueInfo]
     [\pageref{Util-TableQueryValueInfo}]
Query whether or not a specified key is in the table, and optionally
the type and/or number of elements of the value corresponding to this key

\item[Util\_TableQueryMaxKeyLength]
     [\pageref{Util-TableQueryMaxKeyLength}]
Query the maximum key length in a table

\item[Util\_TableQueryNKeys]
     [\pageref{Util-TableQueryNKeys}]
Query the number of key/value entries in a table

\item[Util\_TableSet*]
     [\pageref{Util-TableSet*}]
This is a family of functions, one for each Cactus data type,
to set the value associated with a specified key to be a specified
single (1-element array) value

\item[Util\_TableSet*Array]
     [\pageref{Util-TableSet*Array}]
This is a family of functions, one for each Cactus data type,
to set the value associated with a specified key to be a copy
of a specified array

\item[Util\_TableSetFromString]
     [\pageref{Util-TableSetFromString}]
Sets values in a table based on a string argument (interpreted with
``parameter-file'' semantics)

\item[Util\_TableSetGeneric]
     [\pageref{Util-TableSetGeneric}]
Set the value associated with a specified key to be a specified
single (1-element array) value, whose data type is generic

\item[Util\_TableSetGenericArray]
     [\pageref{Util-TableSetGenericArray}]
Set the value associated with a specified key to be a copy of a
specified array, whose data type is generic

\item[Util\_TableSetString]
     [\pageref{Util-TableSetString}]
Sets the value associated with a specified key in a table, to be
a copy of a specified C-style null-terminated character string

\end{Lentry}

%%%%%%%%%%%%%%%%%%%%%%%%%%%%%%%%%%%%%%%%%%%%%%%%%%%%%%%%%%%%%%%%%%%%%%%%%%%%%%%%

\chapter{Full Descriptions of String Functions}

%%%%%%%%%%%%%%%%%%%%%%%%%%%%%%%%%%%%%%%%

\begin{FunctionDescription}{Util\_StrCmpi}
\label{Util-StrCmpi}
Compare two strings, ignoring upper/lower case.

\begin{SynopsisSection}
\begin{Synopsis}{C}
\begin{verbatim}
#include "util_String.h"
int cmp = Util_StrCmpi(const char *str1, const char *str2);
\end{verbatim}
\end{Synopsis}
\end{SynopsisSection}

\begin{ResultSection}
\begin{Result}{cmp}
An integer which is:\\
\begin{tabular}{@{}rl}
$<0$	& if $\verb|str1| < \verb|str2|$ in lexicographic order
	  ignoring upper/lower case distinctions			\\
$0$	& if $\verb|str1| = \verb|str2|$
	  ignoring upper/lower case distinctions			\\
$>0$	& if $\verb|str1| > \verb|str2|$ in lexicographic order
	  ignoring upper/lower case distinctions			%%%\\
\end{tabular}
\end{Result}
\end{ResultSection}

\begin{ParameterSection}
\begin{Parameter}{str1}
A non-NULL pointer to a (C-style NUL-terminated) string to be compared.
\end{Parameter}
\begin{Parameter}{str2}
A non-NULL pointer to a (C-style NUL-terminated) string to be compared.
\end{Parameter}
\end{ParameterSection}

\begin{Discussion}
The standard C library \verb|strcmp()| function does a {\em case-sensitive\/}
string comparison, \ie{} \verb|strcmp("cactus", "Cactus")| will find the
two strings not equal.  Sometimes it's useful to do {\em case-insensitive\/}
string comparison, where upper/lower case distinctions are ignored.
Many systems provide a \verb|strcasecmp()| or \verb|strcmpi()|
function to do this, but some systems don't, and even on those that do,
the name isn't standardised.  So, Cactus provides its own version,
\verb|Util_StrCmpi()|.

Notice that the return value of \verb|Util_StrCmpi()|, like that of
\verb|strcmp()|, is zero (logical ``false'' in C) for equal strings,
and nonzero (logcla ``true'' in C) for non-equal strings.
Code of the form
\begin{verbatim}
if (Util_StrCmpi(str1, str2))
        { /* strings differ */ }
\end{verbatim}
or
\begin{verbatim}
if (!Util_StrCmpi(str1, str2))
        { /* strings are identical apart from case distinctions */ }
\end{verbatim}
may be confusing to readers, because the sense of the comparison
isn't immediately obvious.  Writing an explicit comparison against zero
make make things clearer:
\begin{verbatim}
if (Util_StrCmpi(str1, str2) != 0)
        { /* strings differ */ }
\end{verbatim}
or
\begin{verbatim}
if (Util_StrCmpi(str1, str2) == 0)
        { /* strings are identical apart from case distinctions */ }
\end{verbatim}

Unfortunately, the basic concept of ``case-insensitive'' string
operations doesn't generalize well to non-English character sets,%%%
\footnote{%%%
	 Hawaiian and Swahili are apparently the only
	 other living languages that use solely the
	 26-letter ``English'' Latin alphabet.
	 }%%%
{} where lower-case $\leftrightarrow$ upper-case mappings may be
context-dependent, many-to-one, and/or time-dependent.%%%
\footnote{%%%
	 For example, the (lower-case) German ``\ss''
	 doesn't have a unique upper-case equivalent:
	 ``\ss'' usually maps to ``SS''
	 (for example ``gro\ss'' $\leftrightarrow$ ``GROSS''),
	 {\em but\/} if that would conflict
	 with another word, then ``\ss'' maps to ``SZ''
	 (for example ``ma\ss{}e'' $\leftrightarrow$ ``MASZE''
	 because there's a different word ``MASSE'').
	 Or at least that's the way it was prior to 1998.
	 The 1998 revisions to German orthography
	 removed the SZ rule, so now (post-1998) the two
	 distinct German words ``masse'' (English ``mass'')
	 and ``ma\ss{}e'' (``measures'') have identical
	 upper-case forms ``MASSE''.  To further complicate
	 matters, (the German-speaking parts of) Switzerland
	 have a slightly different orthography, which never
	 had the SZ rule.

	 French provides another tricky example:
	 In France ``\'e'' $\leftrightarrow$ ``\'E'' and
	 ``\`e'' $\leftrightarrow$ ``\`E'',
	 whereas in (the French-speaking parts of) Canada
	 there are no accents on upper-case letters, so
	 ``\'e'' $\leftrightarrow$ ``E'' and
	 ``\`e'' $\leftrightarrow$ ``E''.
	 }%%%
{}  At present Cactus basically ignores these issues. :(
\end{Discussion}

\begin{SeeAlsoSection}
\begin{SeeAlso}{strcmp()}
Standard C library function (prototype in \verb|<string.h>|)
to compare two strings.
\end{SeeAlso}
\end{SeeAlsoSection}

\begin{ExampleSection}
\begin{Example}{C}
\begin{verbatim}
#include "util_String.h"

/* does the Cactus parameter  driver  specify the PUGH driver? */
/* (Cactus parameters are supposed to be case-insensitive) */
if (Util_StrCmpi(driver, "pugh") == 0)
        { /* PUGH code */ }
else
        { /* non-PUGH code */ }
\end{verbatim}
\end{Example}
\end{ExampleSection}
\end{FunctionDescription}

%%%%%%%%%%%%%%%%%%%%%%%%%%%%%%%%%%%%%%%%

\begin{FunctionDescription}{Util\_Strdup}
\label{Util-Strdup}
``Duplicate'' a string, \ie{} copy it to a newly--\verb|malloc|ed buffer.

\begin{SynopsisSection}
\begin{Synopsis}{C}
\begin{verbatim}
#include "util_String.h"
char* copy = Util_Strdup(const char *str);
\end{verbatim}
\end{Synopsis}
\end{SynopsisSection}

\begin{ResultSection}
\begin{Result}{copy}
A pointer to a buffer obtained from \verb|malloc()|,
which this function sets to a copy of the (C-style NUL-terminated)
string \verb|str|.  This buffer should be freed with \verb|free()|
when it's not needed any more.
\end{Result}
\end{ResultSection}

\begin{ParameterSection}
\begin{Parameter}{str}
A non-NULL pointer to a (C-style NUL-terminated) string.
\end{Parameter}
\end{ParameterSection}

\begin{Discussion}
Many systems have a C library function \verb|strdup|, which
\verb|malloc|s sufficient memory for a copy of its argument string,
does the copy, and returns a pointer to the copy.  However, some
systems lack this function, so Cactus provides its own version,
\verb|Util_Strdup()|.
\end{Discussion}

\begin{SeeAlsoSection}
\begin{SeeAlso}{<stdlib.h>}
System header file containing prototypes for
\verb|malloc()| and \verb|free()|.
\end{SeeAlso}
\begin{SeeAlso}{strcpy()}
Standard C library function (prototype in \verb|<string.h>|)
to copy a string to a buffer.
{\em This does not check that the buffer is big enough to hold the string,
and is thus very dangerous.  Use \verb|Util_Strlcpy()| instead!}
\end{SeeAlso}
\begin{SeeAlso}{Util\_Strlcpy()}
Safely copy a string.
\end{SeeAlso}
\end{SeeAlsoSection}

\begin{ErrorSection}
\begin{Error}{NULL}
\verb|malloc()| was unable to allocate memory for the buffer.
\end{Error}
\end{ErrorSection}

\begin{ExampleSection}
\begin{Example}{C}
\begin{verbatim}
#include "util_String.h"

/*
 * return the (positive) answer to a question,
 * or negative if an error occured
 */
int answer_question(const char* question)
{
/*
 * we need to modify the question string in the process of parsing it
 * but we must not destroy the input ==> copy it and modify the copy
 *
 * ... note the const qualifier on  question_copy  says that
 *     the pointer  question_copy  won't itself change, but
 *     we can modify the string that it points to
 */
char* const question_copy = Util_Strdup(question);
if (question_copy == NULL)
        { return -1; }     /* couldn't get memory for copy buffer */

/* code that will modify  question_copy */

free(question_copy);
return 42;
}
\end{verbatim}
\end{Example}
\end{ExampleSection}
\end{FunctionDescription}

%%%%%%%%%%%%%%%%%%%%%%%%%%%%%%%%%%%%%%%%

\begin{FunctionDescription}{Util\_Strlcat}
\label{Util-Strlcat}
Concatenate strings safely.

\begin{SynopsisSection}
\begin{Synopsis}{C}
\begin{verbatim}
#include "util_String.h"
size_t result_len = Util_Strlcat(char *dst, const char *src, size_t size);
\end{verbatim}
\end{Synopsis}
\end{SynopsisSection}

\begin{ResultSection}
\begin{Result}{result\_len}
The size of the string the function tried to create, i.e.\
the initial \verb|strlen(dst)| plus \verb|strlen(src)|.
\end{Result}
\end{ResultSection}

\begin{ParameterSection}
\begin{Parameter}{dst}
A non-NULL pointer to the (C-style NUL-terminated) destination string.
\end{Parameter}
\begin{Parameter}{src}
A non-NULL pointer to the (C-style NUL-terminated) source string.
\end{Parameter}
\begin{Parameter}{size}
The size of the destination buffer.
\end{Parameter}
\end{ParameterSection}

\begin{Discussion}
The standard \verb|strcat()| and \verb|strcpy()| functions provide
no way to specify the size of the destination buffer, so code using
these functions is often vulnerable to buffer overflows.  The standard
\verb|strncat()| and \verb|strncpy()| functions can be used to write
safe code, but their API is cumbersome, error-prone, and sometimes
surprisingly inefficient:
\begin{itemize}
\item	Their \verb|size| arguments are the number of characters
	{\em remaining\/} in the destination buffer, which must
	often be calculated at run-time, and is prone to off-by-one
	errors.
\item	\verb|strncpy()| doesn't always NUL-terminate the destination string.
\item	\verb|strncpy()| NUL-fills the remainder of the buffer not
	used for the source string; this NUL-filling can be {\em very\/}
	expensive.
\end{itemize}

To solve these problems, the OpenBSD project developed the
\verb|strlcat()| and \verb|strlcpy()| functions.  See
\verb|http://www.openbsd.org/papers/strlcpy-paper.ps|
for a history and general discussion of these functions.
Some other Unix systems (notably Solaris) now provide these,
but many don't, so Cactus provides its own versions,
\verb|Util_Strlcat()| and \verb|Util_Strlcpy()|.

\verb|Util_Strlcat()| appends the NUL-terminated string \verb|src| to the
end of the NUL-terminated string \verb|dst|.  It will append at most
\verb|size - strlen(dst) - 1| characters (hence it never overflows the
destination buffer), and it always leaves \verb|dst| string NUL-terminated.
\end{Discussion}

\begin{SeeAlsoSection}
\begin{SeeAlso}{strcat()}
Standard C library function (prototype in \verb|<string.h>|)
to concatenate two strings.
{\em This does not check that the buffer is big enough to hold the result,
and is thus very dangerous.  Use \verb|Util_Strlcat()| instead!}
\end{SeeAlso}
\begin{SeeAlso}{Util\_Strlcpy()}
Safely copy a string.
\end{SeeAlso}
\end{SeeAlsoSection}

\begin{ExampleSection}
\begin{Example}{C}
\begin{verbatim}
#include "util_String.h"

/*
 * safely concatenate strings s1,s2,s3 into buffer:
 * ... this code is safe (it will never overflow the buffer), but
 *     quick-n-dirty in that it doesn't give any error indication
 *     if the result is truncated to fit in the buffer
 */
#define BUFFER_SIZE     1024
char buffer[BUFFER_SIZE];

Util_Strlcpy(buffer, s1, sizeof(buffer));
Util_Strlcat(buffer, s2, sizeof(buffer));
Util_Strlcat(buffer, s3, sizeof(buffer));
\end{verbatim}
\end{Example}
\begin{Example}{C}
\begin{verbatim}
#include "util_String.h"

#define OK              0
#define ERROR_TRUNC     1

/*
 * safely concatenate strings s1,s2,s3 into buffer[N_buffer];
 * return OK if ok, ERROR_TRUNC if result was truncated to fit in buffer
 */
int cat3(int N_buffer, char buffer[],
         const char s1[], const char s2[], const char s3[])
{
int length;

length = Util_Strlcpy(buffer, s1, N_buffer);
if (length >= N_buffer)
        return ERROR_TRUNC;                   /*** ERROR EXIT ***/

length = Util_Strlcat(buffer, s2, N_buffer);
if (length >= N_buffer)
        return ERROR_TRUNC;                   /*** ERROR EXIT ***/

length = Util_Strlcat(buffer, s3, N_buffer);
if (length >= N_buffer)
        return ERROR_TRUNC;                   /*** ERROR EXIT ***/

return OK;                                    /*** NORMAL RETURN ***/
}
\end{verbatim}
\end{Example}
\end{ExampleSection}
\end{FunctionDescription}

%%%%%%%%%%%%%%%%%%%%%%%%%%%%%%%%%%%%%%%%

\begin{FunctionDescription}{Util\_Strlcpy}
\label{Util-Strlcpy}
Copies a string safely.

\begin{SynopsisSection}
\begin{Synopsis}{C}
\begin{verbatim}
#include "util_String.h"
size_t result_len = Util_Strlcpy(char *dst, const char *src, size_t size);
\end{verbatim}
\end{Synopsis}
\end{SynopsisSection}

\begin{ResultSection}
\begin{Result}{result\_len}
The size of the string the function tried to create, i.e.\
\verb|strlen(src)|.
\end{Result}
\end{ResultSection}

\begin{ParameterSection}
\begin{Parameter}{dst}
A non-NULL pointer to the (C-style NUL-terminated) destination string.
\end{Parameter}
\begin{Parameter}{src}
A non-NULL pointer to the (C-style NUL-terminated) source string.
\end{Parameter}
\begin{Parameter}{size}
The size of the destination buffer.
\end{Parameter}
\end{ParameterSection}

\begin{Discussion}
The standard \verb|strcat()| and \verb|strcpy()| functions provide
no way to specify the size of the destination buffer, so code using
these functions is often vulnerable to buffer overflows.  The standard
\verb|strncat()| and \verb|strncpy()| functions can be used to write
safe code, but their API is cumbersome, error-prone, and sometimes
surprisingly inefficient:
\begin{itemize}
\item	Their \verb|size| arguments are the number of characters
	{\em remaining\/} in the destination buffer, which must
	often be calculated at run-time, and is prone to off-by-one
	errors.
\item	\verb|strncpy()| doesn't always NUL-terminate the destination string.
\item	\verb|strncpy()| NUL-fills the remainder of the buffer not
	used for the source string; this NUL-filling can be {\em very\/}
	expensive.
\end{itemize}

To solve these problems, the OpenBSD project developed the
\verb|strlcat()| and \verb|strlcpy()| functions.  See
\verb|http://www.openbsd.org/papers/strlcpy-paper.ps|
for a history and general discussion of these functions.
Some other Unix systems (notably Solaris) now provide these,
but many don't, so Cactus provides its own versions,
\verb|Util_Strlcat()| and \verb|Util_Strlcpy()|.

\verb|Util_Strlcpy()| copies up to \verb|size-1| characters from the
source string to the destination string, followed by a NUL character
(so \verb|dst| is always NUL-terminated).  Unlike \verb|strncpy()|,
\verb|Util_Strlcpy()| does {\em not\/} fill any left-over space at
the end of the destination buffer with NUL characters.
\end{Discussion}

\begin{SeeAlsoSection}
\begin{SeeAlso}{strcpy()}
Standard C library function (prototype in \verb|<string.h>|)
to copy a string to a buffer.
{\em This does not check that the buffer is big enough to hold the string,
and is thus very dangerous.  Use \verb|Util_Strlcpy()| instead!}
\end{SeeAlso}
\begin{SeeAlso}{Util\_Strdup()}
``Duplicate'' a string, \ie{} copy it to a newly-\verb|malloc|ed buffer.
\end{SeeAlso}
\begin{SeeAlso}{Util\_Strlcat()}
Safely concatenates two strings.
\end{SeeAlso}
\end{SeeAlsoSection}

\begin{ExampleSection}
\begin{Example}{C}
\begin{verbatim}
#include "util_String.h"

/*
 * safely concatenate strings s1,s2,s3 into buffer:
 * ... this code is safe (it will never overflow the buffer), but
 *     quick-n-dirty in that it doesn't give any error indication
 *     if the result is truncated to fit in the buffer
 */
#define BUFFER_SIZE     1024
char buffer[BUFFER_SIZE];

Util_Strlcpy(buffer, s1, sizeof(buffer));
Util_Strlcat(buffer, s2, sizeof(buffer));
Util_Strlcat(buffer, s3, sizeof(buffer));
\end{verbatim}
\end{Example}
\begin{Example}{C}
\begin{verbatim}
#include "util_String.h"

#define OK              0
#define ERROR_TRUNC     1

/*
 * safely concatenate strings s1,s2,s3 into buffer[N_buffer];
 * return OK if ok, ERROR_TRUNC if result was truncated to fit in buffer
 */
int cat3(int N_buffer, char buffer[],
         const char s1[], const char s2[], const char s3[])
{
int length;

length = Util_Strlcpy(buffer, s1, N_buffer);
if (length >= N_buffer)
        return ERROR_TRUNC;                   /*** ERROR EXIT ***/

length = Util_Strlcat(buffer, s2, N_buffer);
if (length >= N_buffer)
        return ERROR_TRUNC;                   /*** ERROR EXIT ***/

length = Util_Strlcat(buffer, s3, N_buffer);
if (length >= N_buffer)
        return ERROR_TRUNC;                   /*** ERROR EXIT ***/

return OK;                                    /*** NORMAL RETURN ***/
}
\end{verbatim}
\end{Example}
\end{ExampleSection}
\end{FunctionDescription}

%%%%%%%%%%%%%%%%%%%%%%%%%%%%%%%%%%%%%%%%%%%%%%%%%%%%%%%%%%%%%%%%%%%%%%%%%%%%%%%%

\chapter{Full Descriptions of Table Functions}

%%%%%%%%%%%%%%%%%%%%%%%%%%%%%%%%%%%%%%%%

\begin{FunctionDescription}{Util\_TableClone}
\label{Util-TableClone}
Creates a new table which is a ``clone'' (exact copy) of an existing
table

\begin{SynopsisSection}
\begin{Synopsis}{C}
\begin{verbatim}
#include "util_ErrorCodes.h"
#include "util_Table.h"
int clone_handle = Util_TableClone(int handle);
\end{verbatim}
\end{Synopsis}
\begin{Synopsis}{Fortran}
\begin{verbatim}
call Util_TableClone(clone_handle, handle)
integer clone_handle, handle
\end{verbatim}
\end{Synopsis}
\end{SynopsisSection}

\begin{ResultSection}
\begin{Result}{clone\_handle ($\ge 0$)}
A handle to the clone table
\end{Result}
\end{ResultSection}

\begin{ParameterSection}
\begin{Parameter}{handle}
Handle to the table to be cloned
\end{Parameter}
\end{ParameterSection}

\begin{Discussion}
Viewing a table as a set of key/value pairs, this function creates
a new table (with the same flags word as the original) containing
copies of all the original table's key/value pairs.  The two tables
are completely independent, \ie{} future changes to one won't affect
the other.

Note that if there are any \verb|CCTK_POINTER| and/or \verb|CCTK_FPOINTER|
values in the table, they are ``shallow copied'', \ie{} the (pointer)
values in the table are copied.  This results in the clone table's
pointer values pointing to the same places as the original table's
pointer values.  Be careful with this!  In particular,
if you're using pointer values in the table to keep track of
\verb|malloc()| memory, be careful not to \verb|free()|
the same block of memory twice!

Note that table iterators are {\em not\/} guaranteed to sequence
through the original and clone tables in the same order.  (This
is a special case of the more general ``non-guarantee'' in
the Section of table iterators in the Users' Guide:
the order of table iterators may differ even between different
tables with identical key/value contents.)
\end{Discussion}

\begin{SeeAlsoSection}
\begin{SeeAlso}{Util\_TableCreate()}
create a table
\end{SeeAlso}
\begin{SeeAlso}{Util\_TableCreateFromString()}
convenience routine to create a table and set key/value entries
in it based on a parameter-file--like character string
\end{SeeAlso}
\begin{SeeAlso}{Util\_TableDestroy()}
destroy a table
\end{SeeAlso}
\end{SeeAlsoSection}

\begin{ErrorSection}
\begin{Error}{UTIL\_ERROR\_NO\_MEMORY}
unable to allocate memory
\end{Error}
\begin{Error}{UTIL\_ERROR\_TABLE\_BAD\_FLAGS}
flags word is negative in the to-be-cloned table
(this indicates an internal error in the table routines,
and should never happen)
\end{Error}
\end{ErrorSection}

\begin{ExampleSection}
\begin{Example}{C}
\begin{verbatim}
#include "util_ErrorCodes.h"
#include "util_Table.h"

/*
 * This function is passed (a handle to) a table containing some entries.
 * It needs to set some additional entries and pass the table to some
 * other function(s), but it also needs to leave the original table
 * intact for other use by the caller.  The solution is to clone the
 * original table and work on the clone, leaving the original table
 * unchanged.
 */
int my_function(int handle, int x, int y)
{
int status;

/* clone the table */
const int clone_handle = Util_TableClone(handle)
if (clone_handle < 0)
        return clone_handle;               /* error in cloning table */

/* now set our entries in the clone table */
status = Util_TableSetInt(clone_handle, x, "x");
if (status < 0)
        return status;                     /* error in setting x */
status = Util_TableSetInt(clone_handle, y, "y");
if (status < 0)
        return status;                     /* error in setting y */

/* ... code to use the clone table ... */
/* ... eg pass clone_handle to other functions ... */

/* we're done with the clone now */
Util_TableDestroy(clone_handle);
return 0;
}
\end{verbatim}
\end{Example}
\end{ExampleSection}
\end{FunctionDescription}

%%%%%%%%%%%%%%%%%%%%%%%%%%%%%%%%%%%%%%%%

\begin{FunctionDescription}{Util\_TableCreate}
\label{Util-TableCreate}
Creates a new (empty) table

\begin{SynopsisSection}
\begin{Synopsis}{C}
\begin{verbatim}
#include "util_ErrorCodes.h"
#include "util_Table.h"
int handle = Util_TableCreate(int flags);
\end{verbatim}
\end{Synopsis}
\begin{Synopsis}{Fortran}
\begin{verbatim}
call Util_TableCreate(handle, flags)
integer handle, flags
\end{verbatim}
\end{Synopsis}
\end{SynopsisSection}

\begin{ResultSection}
\begin{Result}{handle ($\ge 0$)}
A handle to the newly-created table
\end{Result}
\end{ResultSection}

\begin{ParameterSection}
\begin{Parameter}{flags ($\ge 0$)}
A flags word for the table.  This should be the inclusive-or of zero
or more of the \verb|UTIL_TABLE_FLAGS_|* bit masks (defined in
\verb|"util_Table.h"|).  For Fortran users, note that inclusive-or
is the same as sum here, since the bit masks are all disjoint.
\end{Parameter}
\end{ParameterSection}

\begin{Discussion}
We require the flags word to be non-negative so that other
functions can distinguish flags from (negative) error codes.

Any User-defined flag words should use only bit positions at or above\\
\verb|UTIL_TABLE_FLAGS_USER_DEFINED_BASE|, i.e.\ all bit positions
below this are reserved for present of future Cactus use.

At present there is only a single flags-word bit mask defined
in \verb|"util_Table.h"|:
\begin{description}
\item[{\t UTIL\_TABLE\_FLAGS\_CASE\_INSENSITIVE}]\mbox{}\\
        By default keys are treated as C-style character strings,
        and the table functions compare them with the standard C
        \verb|strcmp()| function.
        However, by setting the\\
        \verb|UTIL_TABLE_FLAGS_CASE_INSENSITIVE|
        bit in the flags word, this table's keys may be made
        case-insensitive, \ie{} the table routines then compare
        this table's keys with \verb|Util_StrCmpi()|.%%%
        Note that keys are still {\em stored\/} exactly as
        the caller specifies them (\ie{} they are {\em not\/}
        forced into a canonical case); it's only their
        {\em comparison\/} that's affected by this flag.
\end{description}
\end{Discussion}

\begin{SeeAlsoSection}
\begin{SeeAlso}{Util\_StrCmpi()}
compare two strings, ignoring upper/lower case
\end{SeeAlso}
\begin{SeeAlso}{Util\_TableClone()}
create a new table which is a ``clone'' (exact copy) of an existing
table
\end{SeeAlso}
\begin{SeeAlso}{Util\_TableCreateFromString()}
convenience routine to create a table and set key/value entries
in it based on a parameter-file--like character string
\end{SeeAlso}
\begin{SeeAlso}{Util\_TableDestroy()}
destroy a table
\end{SeeAlso}
\end{SeeAlsoSection}

\begin{ErrorSection}
\begin{Error}{UTIL\_ERROR\_NO\_MEMORY}
unable to allocate memory
\end{Error}
\begin{Error}{UTIL\_ERROR\_TABLE\_BAD\_FLAGS}
flags word is negative
\end{Error}
\end{ErrorSection}

\begin{ExampleSection}
\begin{Example}{C}
\begin{verbatim}
#include "util_ErrorCodes.h"
#include "util_Table.h"

/* create a table, simplest case */
int handle = Util_TableCreate(UTIL_TABLE_FLAGS_DEFAULT);

/* create a table whose keys will be treated as case-insensitive */
int handle2 = Util_TableCreate(UTIL_TABLE_FLAGS_CASE_INSENSITIVE);
\end{verbatim}
\end{Example}
\end{ExampleSection}
\end{FunctionDescription}

%%%%%%%%%%%%%%%%%%%%%%%%%%%%%%%%%%%%%%%%

\begin{FunctionDescription}{Util\_TableCreateFromString}
\label{Util-TableCreateFromString}
Creates a new table (with the case-insensitive flag set) and sets
values in it based on a string argument (interpreted with
``parameter-file'' semantics)

\begin{SynopsisSection}
\begin{Synopsis}{C}
\begin{verbatim}
#include "util_ErrorCodes.h"
#include "util_Table.h"
int handle = Util_TableCreateFromString(const char *string);
\end{verbatim}
\end{Synopsis}
\begin{Synopsis}{Fortran}
\begin{verbatim}
call Util_TableCreateFromString(handle, string)
integer       handle
character*(*) string
\end{verbatim}
\end{Synopsis}
\end{SynopsisSection}

\begin{ResultSection}
\begin{Result}{handle ($\ge 0$)}
a handle to the newly-created table
\end{Result}
\end{ResultSection}

\begin{ParameterSection}
\begin{Parameter}{string}
a pointer to a C-style null-terminated string specifying the table
contents; see the description for \verb|Util_TableSetFromString()|
for a full description of the syntax and semantics of this string
\end{Parameter}
\end{ParameterSection}

\begin{SeeAlsoSection}
\begin{SeeAlso}{Util\_TableClone()}
Create a new table which is a ``clone'' (exact copy) of an existing
table
\end{SeeAlso}
\begin{SeeAlso}{Util\_TableCreate()}
create a table
\end{SeeAlso}
\begin{SeeAlso}{Util\_TableSetFromString()}
sets values in a table based on a string argument
\end{SeeAlso}
\begin{SeeAlso}{Util\_TableSetInt()}
store a \verb|CCTK_INT| value in a table
\end{SeeAlso}
\begin{SeeAlso}{Util\_TableSetReal()}
store a \verb|CCTK_REAL| value in a table
\end{SeeAlso}
\end{SeeAlsoSection}

\begin{ErrorSection}
\begin{Error}{UTIL\_ERROR\_NO\_MEMORY}
unable to allocate memory
\end{Error}
\begin{Error}{UTIL\_ERROR\_BAD\_KEY}
invalid input: key contains invalid character
\end{Error}
\begin{Error}{UTIL\_ERROR\_BAD\_INPUT}
invalid input: can't parse input string
\end{Error}
\begin{Error}{\rm other error codes}
this function may also return any error codes returned by
\verb|Util_TableCreate()| or \verb|Util_TableSetFromString()|
\end{Error}
\end{ErrorSection}

\begin{ExampleSection}
\begin{Example}{C}
\begin{verbatim}
#include "util_ErrorCodes.h"
#include "util_Table.h"

int handle = Util_TableCreateFromString("order = 3\t"
                                        "myreal = 42.314159\t"
                                        "mystring = 'hello'\t"
                                        "myarray = { 0 1 2 3 }");

/* equivalent code to the above */
int handle = Util_TableCreate(UTIL_TABLE_FLAGS_CASE_INSENSITIVE);
Util_TableSetFromString(handle, "order = 3\t"
                                "myreal = 42.314159\t"
                                "mystring = 'hello'"
                                "myarray = { 0 1 2 3 }");

/* also equivalent to the above */
int handle = Util_TableCreate(UTIL_TABLE_FLAGS_CASE_INSENSITIVE);
CCTK_INT array[] = {0, 1, 2, 3};

Util_TableSetInt(handle, 3, "order");
Util_TableSetReal(handle, 42.314159, "myreal");
Util_TableSetString(handle, "hello", "mystring");
Util_TableSetIntArray(handle, 4, array, "myarray");
\end{verbatim}
\end{Example}
\end{ExampleSection}
\end{FunctionDescription}

%%%%%%%%%%%%%%%%%%%%%%%%%%%%%%%%%%%%%%%%

\begin{FunctionDescription}{Util\_TableDeleteKey}
\label{Util-TableDeleteKey}
Deletes a specified key/value entry from a table

\begin{SynopsisSection}
\begin{Synopsis}{C}
\begin{verbatim}
#include "util_ErrorCodes.h"
#include "util_Table.h"
int key_exists = Util_TableDeleteKey(int handle, const char *key);
\end{verbatim}
\end{Synopsis}
\begin{Synopsis}{Fortran}
\begin{verbatim}
call Util_TableDeleteKey(key_exists, handle, key)
integer        key_exists, handle
character*(*)  key
\end{verbatim}
\end{Synopsis}
\end{SynopsisSection}

\begin{ResultSection}
\begin{Result}{\rm 0}
ok (key existed before this call, and has now been deleted)
\end{Result}
\end{ResultSection}

\begin{ParameterSection}
\begin{Parameter}{handle ($\ge 0$)}
handle to the table
\end{Parameter}
\begin{Parameter}{key}
a pointer to the key (a C-style null-terminated string)
\end{Parameter}
\end{ParameterSection}

\begin{Discussion}
This function invalidates any iterators for the table which are
not in the ``null-pointer'' state.
\end{Discussion}

\begin{ErrorSection}
\begin{Error}{UTIL\_ERROR\_BAD\_HANDLE}
handle is invalid
\end{Error}
\begin{Error}{UTIL\_ERROR\_TABLE\_BAD\_KEY}
key contains '/' character
\end{Error}
\begin{Error}{UTIL\_ERROR\_TABLE\_NO\_SUCH\_KEY}
no such key in table
\end{Error}
\end{ErrorSection}
\end{FunctionDescription}

%%%%%%%%%%%%%%%%%%%%%%%%%%%%%%%%%%%%%%%%

\begin{FunctionDescription}{Util\_TableDestroy}
\label{Util-TableDestroy}
Destroys a table

\begin{SynopsisSection}
\begin{Synopsis}{C}
\begin{verbatim}
#include "util_ErrorCodes.h"
#include "util_Table.h"
int status = Util_TableDestroy(int handle);
\end{verbatim}
\end{Synopsis}
\begin{Synopsis}{Fortran}
\begin{verbatim}
call Util_TableDestroy(status, handle)
integer status, handle
\end{verbatim}
\end{Synopsis}
\end{SynopsisSection}

\begin{ResultSection}
\begin{Result}{\rm 0}
ok
\end{Result}
\end{ResultSection}

\begin{ParameterSection}
\begin{Parameter}{handle ($\ge 0$)}
handle to the table
\end{Parameter}
\end{ParameterSection}

\begin{Discussion}
Of course, this function invalidates any and all iterators for the table. :)
\end{Discussion}

\begin{SeeAlsoSection}
\begin{SeeAlso}{Util\_TableClone()}
Create a new table which is a ``clone'' (exact copy) of an existing
table
\end{SeeAlso}
\begin{SeeAlso}{Util\_TableCreate()}
create a table
\end{SeeAlso}
\begin{SeeAlso}{Util\_TableCreateFromString()}
convenience routine to create a table and set key/value entries
in it based on a parameter-file--like character string
\end{SeeAlso}
\end{SeeAlsoSection}

\begin{ErrorSection}
\begin{Error}{UTIL\_ERROR\_BAD\_HANDLE}
handle is invalid
\end{Error}
\end{ErrorSection}

\begin{ExampleSection}
\begin{Example}{C}
\begin{verbatim}
#include "util_ErrorCodes.h"
#include "util_Table.h"

/* create a table */
int handle = Util_TableCreate(UTIL_TABLE_FLAGS_DEFAULT);

/* do things with the table: put values in it, */
/* pass its handle to other functions, etc etc */
/* ... */

/* at this point we (and all other functions we */
/* may call in the future) are done with the table */
Util_TableDestroy(handle);
\end{verbatim}
\end{Example}
\end{ExampleSection}
\end{FunctionDescription}

%%%%%%%%%%%%%%%%%%%%%%%%%%%%%%%%%%%%%%%%

\begin{FunctionDescription}{Util\_TableGet*}
\label{Util-TableGet*}
This is a family of functions, one for each Cactus data type,
to get the single (1-element array) value, or more generally the
first array element of the value, associated with a specified key
in a key/value table.

\begin{SynopsisSection}
\begin{Synopsis}{C}
\begin{verbatim}
#include "util_ErrorCodes.h"
#include "util_Table.h"
int N_elements = Util_TableGetXxx(int handle,
                                  CCTK_XXX *value,
                                  const char *key);
\end{verbatim}
where \verb|XXX| is one of
   \verb|POINTER|, \verb|FPOINTER|%%%
\footnote{%%%
	 For backwards compatability the function
	 {\tt Util\_TableGetFnPointer()} is also provided
	 as an alias for {\tt Util\_TableGetFPointer()}.
	 This is deprecated as of Cactus 4.0 beta 13.
	 }%%%
,
   \verb|CHAR|,
   \verb|INT|, \verb|INT2|, \verb|INT4|, \verb|INT8|,
   \verb|REAL|, \verb|REAL4|, \verb|REAL8|, \verb|REAL16|,
   \verb|COMPLEX|, \verb|COMPLEX8|, \verb|COMPLEX16|, \verb|COMPLEX32|
(not all of these may be supported on any given system)
\end{Synopsis}
\begin{Synopsis}{Fortran}
\begin{verbatim}
call Util_TableGetXxx(N_elements, handle, value, key)
integer        N_elements, handle
CCTK_XXX       value
character*(*)  key
\end{verbatim}
\end{Synopsis}
\end{SynopsisSection}

\begin{ResultSection}
\begin{Result}{N\_elements}
the number of array elements in the value
\end{Result}
\end{ResultSection}

\begin{ParameterSection}
\begin{Parameter}{handle ($\ge 0$)}
handle to the table
\end{Parameter}
\begin{Parameter}{value}
a pointer to where this function should store a copy of the value
(or more generally the first array element of the value) associated
with the specified key,
or NULL pointer to skip storing this
\end{Parameter}
\begin{Parameter}{key}
a pointer to the key (a C-style null-terminated string)
\end{Parameter}
\end{ParameterSection}

\begin{Discussion}
Note that it is {\em not\/} an error for the value to actually have
$> 1$ array elements; in this case only the first element is stored.
The rationale for this design is that the caller may know or suspect
that the value is a large array, but may only want the first array
element; in this case this design avoids the caller having to allocate
a large buffer unnecessarily.

In contrast, it {\em is\/} an error for the value to actually be an
empty (0-length) array, because then there is no ``first array element''
to get.

It is also an error for the value to actually have a different type
than \verb|CCTK_XXX|.

If any error code is returned, the user's value buffer
(pointed to by \verb|value| if this is non-NULL) is unchanged.
\end{Discussion}

\begin{SeeAlsoSection}
\begin{SeeAlso}{Util\_TableCreateFromString()}
convenience routine to create a table and set key/value entries
in it based on a parameter-file--like character string
\end{SeeAlso}
\begin{SeeAlso}{Util\_TableGet*Array()}
get an array value
\end{SeeAlso}
\begin{SeeAlso}{Util\_TableGet*String()}
get a character-string value
\end{SeeAlso}
\begin{SeeAlso}{Util\_TableSet*()}
set a single (1-element array) value
\end{SeeAlso}
\begin{SeeAlso}{Util\_TableSet*Array()}
set an array value
\end{SeeAlso}
\begin{SeeAlso}{Util\_TableSetGeneric()}
set a single (1-element array) value with generic data type
\end{SeeAlso}
\begin{SeeAlso}{Util\_TableSetGenericArray()}
set an array value with generic data type
\end{SeeAlso}
\begin{SeeAlso}{Util\_TableSetFromString()}
convenience routine to set key/value entries in a table based on a
parameter-file--like character string
\end{SeeAlso}
\begin{SeeAlso}{Util\_TableSetString()}
set a character-string value
\end{SeeAlso}
\end{SeeAlsoSection}

\begin{ErrorSection}
\begin{Error}{UTIL\_ERROR\_BAD\_HANDLE}
handle is invalid
\end{Error}
\begin{Error}{UTIL\_ERROR\_TABLE\_BAD\_KEY}
key contains '/' character
\end{Error}
\begin{Error}{UTIL\_ERROR\_TABLE\_NO\_SUCH\_KEY}
no such key in table
\end{Error}
\begin{Error}{UTIL\_ERROR\_TABLE\_WRONG\_DATA\_TYPE}
value has data type other than \verb|CCTK_TYPE|
\end{Error}
\begin{Error}{UTIL\_ERROR\_TABLE\_VALUE\_IS\_EMPTY}
value is an empty (0-element) array
\end{Error}
\end{ErrorSection}

\begin{ExampleSection}
\begin{Example}{C}
\begin{verbatim}
#include "util_ErrorCodes.h"
#include "util_Table.h"

#define N_DIGITS        5
static const CCTK_INT pi_digits[N_DIGITS] = {3, 14, 159, 2653, 58979};

int N;
CCTK_INT x;
int handle = Util_TableCreate(UTIL_TABLE_FLAGS_DEFAULT);

Util_TableSetIntArray(handle, N_DIGITS, pi_digits, "digits of pi");
Util_TableSetIntArray(handle, 0, pi_digits, "empty array");

/* gets N = 5, x = 3 */
N = Util_TableGetInt(handle, &x, "digits of pi");

/* gets N = UTIL_ERROR_TABLE_VALUE_IS_EMPTY */
N = Util_TableGetInt(handle, &x, "empty array");
\end{verbatim}
\end{Example}
\end{ExampleSection}
\end{FunctionDescription}

%%%%%%%%%%%%%%%%%%%%%%%%%%%%%%%%%%%%%%%%

\begin{FunctionDescription}{Util\_TableGet*Array}
\label{Util-TableGet*Array}
This is a family of functions, one for each Cactus data type,
to get a copy of the value associated with a specified key, and store
it (more accurately, as much of it as will fit) in a specified array

\begin{SynopsisSection}
\begin{Synopsis}{C}
\begin{verbatim}
#include "util_ErrorCodes.h"
#include "util_Table.h"
int N_elements = Util_TableGetXxxArray(int handle,
                                       int N_array, CCTK_XXX array[],
                                       const char *key);
\end{verbatim}
where \verb|XXX| is one of
   \verb|POINTER|, \verb|FPOINTER|%%%
\footnote{%%%
	 For backwards compatability the function
	 {\tt Util\_TableGetFnPointerArray()} is also provided
	 as an alias for {\tt Util\_TableGetFPointerArray()}.
	 This is deprecated as of Cactus 4.0 beta 13.
	 }%%%
,
   \verb|CHAR|,
   \verb|INT|, \verb|INT2|, \verb|INT4|, \verb|INT8|,
   \verb|REAL|, \verb|REAL4|, \verb|REAL8|, \verb|REAL16|,
   \verb|COMPLEX|, \verb|COMPLEX8|, \verb|COMPLEX16|, \verb|COMPLEX32|
(not all of these may be supported on any given system)
\end{Synopsis}
\begin{Synopsis}{Fortran}
\begin{verbatim}
call Util_TableGetXxxArray(N_elements, handle, N_array, array, key)
integer        N_elements, handle, N_array
CCTK_XXX(*)    array
character*(*)  key
\end{verbatim}
\end{Synopsis}
\end{SynopsisSection}

\begin{ResultSection}
\begin{Result}{N\_elements}
the number of array elements in the value
\end{Result}
\end{ResultSection}

\begin{ParameterSection}
\begin{Parameter}{handle ($\ge 0$)}
handle to the table
\end{Parameter}
\begin{Parameter}{N\_array}
the number of array elements in \verb|array[]|
(must be $\ge 0$ if \verb|array != NULL|)
\end{Parameter}
\begin{Parameter}{array}
a pointer to where this function should store (up to \verb|N_array|
elements of) a copy of the value associated with the specified key,
or NULL pointer to skip storing this
\end{Parameter}
\begin{Parameter}{key}
a pointer to the key (a C-style null-terminated string)
\end{Parameter}
\end{ParameterSection}

\begin{Discussion}
Note that it is {\em not\/} an error for the value to actually have
$> \verb|N_array|$ array elements; in this case only the first \verb|N_array|
elements are stored.  The caller can detect this by comparing the
return value with \verb|N_array|.
The rationale for this design is that the caller may know or suspect
that the value is a large array, but may only want the first few array
elements; in this case this design avoids the caller having to allocate
a large buffer unnecessarily.

It is also {\em not\/} an error for the value to actually have
$< \verb|N_array|$ array elements; again the caller can detect this by
comparing the return value with \verb|N_array|.

It {\em is\/} an error for the value to actually have a different type
than \verb|CCTK_XXX|.

If any error code is returned, the user's value buffer
(pointed to by \verb|array| if this is non-NULL) is unchanged.
\end{Discussion}

\begin{SeeAlsoSection}
\begin{SeeAlso}{Util\_TableCreateFromString()}
convenience routine to create a table and set key/value entries
in it based on a parameter-file--like character string
\end{SeeAlso}
\begin{SeeAlso}{Util\_TableGet*()}
get a single (1-element array) value,
or more generally the first array element of an array value
\end{SeeAlso}
\begin{SeeAlso}{Util\_TableGetGeneric()}
get a single (1-element array) value with generic data type
\end{SeeAlso}
\begin{SeeAlso}{Util\_TableGetGenericArray()}
get an array value with generic data type
\end{SeeAlso}
\begin{SeeAlso}{Util\_TableGet*String()}
get a character-string value
\end{SeeAlso}
\begin{SeeAlso}{Util\_TableSet*()}
set a single (1-element array) value
\end{SeeAlso}
\begin{SeeAlso}{Util\_TableSet*Array()}
set an array value
\end{SeeAlso}
\begin{SeeAlso}{Util\_TableSetGeneric()}
set a single (1-element array) value with generic data type
\end{SeeAlso}
\begin{SeeAlso}{Util\_TableSetGenericArray()}
set an array value with generic data type
\end{SeeAlso}
\begin{SeeAlso}{Util\_TableSetFromString()}
convenience routine to set key/value entries in a table based on a
parameter-file--like character string
\end{SeeAlso}
\begin{SeeAlso}{Util\_TableSetString()}
set a character-string value
\end{SeeAlso}
\end{SeeAlsoSection}

\begin{ErrorSection}
\begin{Error}{UTIL\_ERROR\_BAD\_HANDLE}
handle is invalid
\end{Error}
\begin{Error}{UTIL\_ERROR\_TABLE\_BAD\_KEY}
key contains '/' character
\end{Error}
\begin{Error}{UTIL\_ERROR\_BAD\_INPUT}
\verb|array != NULL| and \verb|N_array| $< 0$
\end{Error}
\begin{Error}{UTIL\_ERROR\_TABLE\_NO\_SUCH\_KEY}
no such key in table
\end{Error}
\begin{Error}{UTIL\_ERROR\_TABLE\_WRONG\_DATA\_TYPE}
value has data type other than \verb|CCTK_TYPE|
\end{Error}
\end{ErrorSection}

\begin{ExampleSection}
\begin{Example}{C}
\begin{verbatim}
#include "util_ErrorCodes.h"
#include "util_Table.h"

#define N_STUFF         3
static const CCTK_REAL stuff[N_STUFF] = {42.0, 69.0, 105.5};

#define N_OUTPUT        2
CCTK_INT output[N_OUTPUT];

int N;
int handle = Util_TableCreate(UTIL_TABLE_FLAGS_DEFAULT);

Util_TableSetRealArray(handle, N_STUFF, stuff, "blah blah blah");

/* gets N = 3, output[0] = 42.0, output[1] = 69.0 */
N = Util_TableGetRealArray(handle, N_OUTPUT, output, "blah blah blah");
\end{verbatim}
\end{Example}
\end{ExampleSection}
\end{FunctionDescription}

%%%%%%%%%%%%%%%%%%%%%%%%%%%%%%%%%%%%%%%%

\begin{FunctionDescription}{Util\_TableGetGeneric}
\label{Util-TableGetGeneric}
Get the single (1-element array) value, or more generally the
first array element of the value, associated with a specified key
in a key/value table; the value's data type is generic.  That
is, the value is specified by a \verb|CCTK_VARIABLE_|* type code
and a \verb|void *| pointer.

\begin{SynopsisSection}
\begin{Synopsis}{C}
\begin{verbatim}
#include "util_ErrorCodes.h"
#include "util_Table.h"
int N_elements = Util_TableGetGeneric(int handle,
                                      int type_code,
                                      void *value,
                                      const char *key);
\end{verbatim}
\end{Synopsis}
\begin{Synopsis}{Fortran}
\begin{verbatim}
call Util_TableGetGeneric(N_elements, handle, type_code, value, key)
integer        N_elements, handle, type_code
CCTK_POINTER   value
character*(*)  key
\end{verbatim}
\end{Synopsis}
\end{SynopsisSection}

\begin{ResultSection}
\begin{Result}{N\_elements}
the number of array elements in the value
\end{Result}
\end{ResultSection}

\begin{ParameterSection}
\begin{Parameter}{handle ($\ge 0$)}
handle to the table
\end{Parameter}
\begin{Parameter}{type\_code}
the value's type code
(one of the \verb|CCTK_VARIABLE_|* constants from \verb|"cctk_Constants.h"|)
\end{Parameter}
\begin{Parameter}{value}
a pointer to where this function should store a copy of the value
(or more generally the first array element of the value) associated
with the specified key,
or NULL pointer to skip storing this
\end{Parameter}
\begin{Parameter}{key}
a pointer to the key (a C-style null-terminated string)
\end{Parameter}
\end{ParameterSection}

\begin{Discussion}
Note that it is {\em not\/} an error for the value to actually have
$> 1$ array elements; in this case only the first element is stored.
The rationale for this design is that the caller may know or suspect
that the value is a large array, but may only want the first array
element; in this case this design avoids the caller having to allocate
a large buffer unnecessarily.

In contrast, it {\em is\/} an error for the value to actually be an
empty (0-length) array, because then there is no ``first array element''
to get.

It is also an error for the value to actually have a different type
than that specified by \verb|type_code|.

If any error code is returned, the user's value buffer
(pointed to by \verb|value| if this is non-NULL) is unchanged.
\end{Discussion}

\begin{SeeAlsoSection}
\begin{SeeAlso}{Util\_TableCreateFromString()}
convenience routine to create a table and set key/value entries
in it based on a parameter-file--like character string
\end{SeeAlso}
\begin{SeeAlso}{Util\_TableGet*()}
get a single (1-element array) value
\end{SeeAlso}
\begin{SeeAlso}{Util\_TableGet*Array()}
get an array value
\end{SeeAlso}
\begin{SeeAlso}{Util\_TableGet*String()}
get a character-string value
\end{SeeAlso}
\begin{SeeAlso}{Util\_TableItQueryValueInfo()}
query key present/absent in table, and optionally type and/or number
of elements, using an iterator
\end{SeeAlso}
\begin{SeeAlso}{Util\_TableQueryValueInfo()}
query key present/absent in table, and optionally type and/or number
of elements
\end{SeeAlso}
\begin{SeeAlso}{Util\_TableSet*()}
set a single (1-element array) value
\end{SeeAlso}
\begin{SeeAlso}{Util\_TableSet*Array()}
set an array value
\end{SeeAlso}
\begin{SeeAlso}{Util\_TableSetGeneric()}
set a single (1-element array) value with generic data type
\end{SeeAlso}
\begin{SeeAlso}{Util\_TableSetGenericArray()}
set an array value with generic data type
\end{SeeAlso}
\begin{SeeAlso}{Util\_TableSetFromString()}
convenience routine to set key/value entries in a table based on a
parameter-file--like character string
\end{SeeAlso}
\begin{SeeAlso}{Util\_TableSetString()}
set a character-string value
\end{SeeAlso}
\end{SeeAlsoSection}

\begin{ErrorSection}
\begin{Error}{UTIL\_ERROR\_BAD\_HANDLE}
handle is invalid
\end{Error}
\begin{Error}{UTIL\_ERROR\_TABLE\_BAD\_KEY}
key contains '/' character
\end{Error}
\begin{Error}{UTIL\_ERROR\_TABLE\_NO\_SUCH\_KEY}
no such key in table
\end{Error}
\begin{Error}{UTIL\_ERROR\_TABLE\_WRONG\_DATA\_TYPE}
value has data type other than \verb|CCTK_TYPE|
\end{Error}
\begin{Error}{UTIL\_ERROR\_TABLE\_VALUE\_IS\_EMPTY}
value is an empty (0-element) array
\end{Error}
\end{ErrorSection}

\begin{ExampleSection}
\begin{Example}{C}
\begin{verbatim}
#include "util_ErrorCodes.h"
#include "util_Table.h"
#include "cctk_Constants.h"

#define N_DIGITS        5
static const CCTK_INT pi_digits[N_DIGITS] = {3, 14, 159, 2653, 58979};

int N;
CCTK_INT x;
void *xptr = (void *) &x;
int handle = Util_TableCreate(UTIL_TABLE_FLAGS_DEFAULT);

Util_TableSetIntArray(handle, N_DIGITS, pi_digits, "digits of pi");
Util_TableSetIntArray(handle, 0, pi_digits, "empty array");

/* gets N = 5, x = 3 */
N = Util_TableGetGeneric(handle, CCTK_VARIABLE_INT, &x, "the answer");

/* gets N = UTIL_ERROR_TABLE_VALUE_IS_EMPTY, leaves x unchanged */
N = Util_TableGetGeneric(handle, CCTK_VARIABLE_INT, &x, "empty array");
\end{verbatim}
\end{Example}
\end{ExampleSection}
\end{FunctionDescription}

%%%%%%%%%%%%%%%%%%%%%%%%%%%%%%%%%%%%%%%%

\begin{FunctionDescription}{Util\_TableGetGenericArray}
\label{Util-TableGetGenericArray}
Get a copy of the value associated with a specified key, and store
it (more accurately, as much of it as will fit) in a specified array;
the array's data type is generic.  That is the array is specified by
a \verb|CCTK_VARIABLE_|* type code, a count of the number of array
elements, and a \verb|void *| pointer.

\begin{SynopsisSection}
\begin{Synopsis}{C}
\begin{verbatim}
#include "util_ErrorCodes.h"
#include "util_Table.h"
int N_elements = Util_TableGetGenericArray(int handle,
                                           int type_code,
                                           int N_array, void *array,
                                           const char *key);
\end{verbatim}
\end{Synopsis}
\begin{Synopsis}{Fortran}
\begin{verbatim}
call Util_TableGetGenericArray(N_elements,
.                              handle,
.                              type_code,
.                              N_array, array,
.                              key)
integer        N_elements, handle, type_code, N_array
CCTK_POINTER   array
character*(*)  key
\end{verbatim}
\end{Synopsis}
\end{SynopsisSection}

\begin{ResultSection}
\begin{Result}{N\_elements}
the number of array elements in the value
\end{Result}
\end{ResultSection}

\begin{ParameterSection}
\begin{Parameter}{handle ($\ge 0$)}
handle to the table
\end{Parameter}
\begin{Parameter}{type\_code}
the value's type code
(one of the \verb|CCTK_VARIABLE_|* constants from \verb|"cctk_Constants.h"|)
\end{Parameter}
\begin{Parameter}{N\_array}
the number of array elements in \verb|array[]|
(must be $\ge 0$ if \verb|array != NULL|)
\end{Parameter}
\begin{Parameter}{array}
a pointer to where this function should store (up to \verb|N_array|
elements of) a copy of the value associated with the specified key,
or NULL pointer to skip storing this
\end{Parameter}
\begin{Parameter}{key}
a pointer to the key (a C-style null-terminated string)
\end{Parameter}
\end{ParameterSection}

\begin{Discussion}
Note that it is {\em not\/} an error for the value to actually have
$> \verb|N_array|$ array elements; in this case only the first \verb|N_array|
elements are stored.  The caller can detect this by comparing the
return value with \verb|N_array|.
The rationale for this design is that the caller may know or suspect
that the value is a large array, but may only want the first few array
elements; in this case this design avoids the caller having to allocate
a large buffer unnecessarily.

It is also {\em not\/} an error for the value to actually have
$< \verb|N_array|$ array elements; again the caller can detect this by
comparing the return value with \verb|N_array|.

It {\em is\/} an error for the value to actually have a different type
than that specified by \verb|type_code|.

If any error code is returned, the user's value buffer
(pointed to by \verb|array| if this is non-NULL) is unchanged.
\end{Discussion}

\begin{SeeAlsoSection}
\begin{SeeAlso}{Util\_TableCreateFromString()}
convenience routine to create a table and set key/value entries
in it based on a parameter-file--like character string
\end{SeeAlso}
\begin{SeeAlso}{Util\_TableGet*()}
get a single (1-element array) value,
or more generally the first array element of an array value
\end{SeeAlso}
\begin{SeeAlso}{Util\_TableGetGeneric()}
get a single (1-element array) value with generic data type
\end{SeeAlso}
\begin{SeeAlso}{Util\_TableGetGenericArray()}
get an array value with generic data type
\end{SeeAlso}
\begin{SeeAlso}{Util\_TableGet*String()}
get a character-string value
\end{SeeAlso}
\begin{SeeAlso}{Util\_TableItQueryValueInfo()}
query key present/absent in table, and optionally type and/or number
of elements, using an iterator
\end{SeeAlso}
\begin{SeeAlso}{Util\_TableQueryValueInfo()}
query key present/absent in table, and optionally type and/or number
of elements
\end{SeeAlso}
\begin{SeeAlso}{Util\_TableSet*()}
set a single (1-element array) value
\end{SeeAlso}
\begin{SeeAlso}{Util\_TableSet*Array()}
set an array value
\end{SeeAlso}
\begin{SeeAlso}{Util\_TableSetGeneric()}
set a single (1-element array) value with generic data type
\end{SeeAlso}
\begin{SeeAlso}{Util\_TableSetGenericArray()}
set an array value with generic data type
\end{SeeAlso}
\begin{SeeAlso}{Util\_TableSetFromString()}
convenience routine to set key/value entries in a table based on a
parameter-file--like character string
\end{SeeAlso}
\begin{SeeAlso}{Util\_TableSetString()}
set a character-string value
\end{SeeAlso}
\end{SeeAlsoSection}

\begin{ErrorSection}
\begin{Error}{UTIL\_ERROR\_BAD\_HANDLE}
handle is invalid
\end{Error}
\begin{Error}{UTIL\_ERROR\_TABLE\_BAD\_KEY}
key contains '/' character
\end{Error}
\begin{Error}{UTIL\_ERROR\_BAD\_INPUT}
\verb|array != NULL| and \verb|N_array| $< 0$
\end{Error}
\begin{Error}{UTIL\_ERROR\_TABLE\_NO\_SUCH\_KEY}
no such key in table
\end{Error}
\begin{Error}{UTIL\_ERROR\_TABLE\_WRONG\_DATA\_TYPE}
value has data type other than \verb|CCTK_TYPE|
\end{Error}
\end{ErrorSection}

\begin{ExampleSection}
\begin{Example}{C}
\begin{verbatim}
#include "util_ErrorCodes.h"
#include "util_Table.h"

#define N_STUFF         3
static const CCTK_REAL stuff[N_STUFF] = {42.0, 69.0, 105.5};

#define N_OUTPUT        2
CCTK_INT output[N_OUTPUT];

int N;
int handle = Util_TableCreate(UTIL_TABLE_FLAGS_DEFAULT);

Util_TableSetRealArray(handle, N_STUFF, stuff, "stuff");

/* gets N = UTIL_ERROR_TABLE_WRONG_DATA_TYPE, output[] unchanged */
N = Util_TableGetGenericArray(handle,
                              CCTK_VARIABLE_INT,
                              N_OUTPUT, output,
                              "stuff");
/* gets N = 3, output[0] = 42.0, output[1] = 69.0 */
N = Util_TableGetGenericArray(handle,
                              CCTK_VARIABLE_REAL,
                              N_OUTPUT, output,
                              "stuff");
\end{verbatim}
\end{Example}
\end{ExampleSection}
\end{FunctionDescription}

%%%%%%%%%%%%%%%%%%%%%%%%%%%%%%%%%%%%%%%%

\begin{FunctionDescription}{Util\_TableGetString}
\label{Util-TableGetString}
Gets a copy of the character-string value associated with a specified
key in a table, and stores it (more accurately, as much of it as will fit)
in a specified character string

\begin{SynopsisSection}
\begin{Synopsis}{C}
\begin{verbatim}
#include "util_ErrorCodes.h"
#include "util_Table.h"
int length = Util_TableGetString(int handle,
                                 int buffer_length, char buffer[],
                                 const char *key);
\end{verbatim}
\end{Synopsis}
\end{SynopsisSection}

\begin{ResultSection}
\begin{ResultNote}
Results are the same as all the other \verb|Util_TableGet|* functions:
\end{ResultNote}
\begin{Result}{length}
the length of the string
(C \verb|strlen()| semantics, \ie{} {\em not\/} including
the terminating null character)
\end{Result}
\end{ResultSection}

\begin{ParameterSection}
\begin{Parameter}{handle ($\ge 0$)}
handle to the table
\end{Parameter}
\begin{Parameter}{buffer\_length}
the length (\verb|sizeof()|) of \verb|buffer[]|
(must be $\ge 1$ if \verb|buffer != NULL|)
\end{Parameter}
\begin{Parameter}{buffer}
a pointer to a buffer into which this function should store
(at most \verb|buffer_length-1|  characters of) the value,
terminated by a null character as usual for C strings,
or NULL pointer to skip storing this
\end{Parameter}
\begin{Parameter}{key}
a pointer to the key (a C-style null-terminated string)
\end{Parameter}
\end{ParameterSection}

\begin{Discussion}
This function assumes that the string is stored as an array of
\verb|CCTK_CHAR|s, {\em not\/} including a terminating null character.

This function differs from \verb|Util_TableGetCharArray()| in two ways:
It explicitly provides a terminating null character for C-style strings,
and it explicitly checks for the string being too long to fit in the buffer
(in which case it returns \verb|UTIL_ERROR_TABLE_STRING_TRUNCATED|).

If the error code \verb|UTIL_ERROR_TABLE_STRING_TRUNCATED| is returned,
then the first \verb|buffer_length-1| characters of the string are
returned in the user's buffer (assuming \verb|buffer| is non-NULL),
followed by a null character to properly terminate the string in the
buffer.  If any other error code is returned, the user's value buffer
(pointed to by \verb|buffer| if this is non-NULL) is unchanged.

To find out how long the string is (and thus how big of a buffer you
need to allocate to avoid having the string truncated), you can call
this function with $\verb|buffer_length| = 0$ and $\verb|buffer| = \verb|NULL|$
(or actually anything you want); the return result will give the
string length.
\end{Discussion}

\begin{SeeAlsoSection}
\begin{SeeAlso}{Util\_TableCreateFromString()}
convenience routine to create a table and set key/value entries
in it based on a parameter-file--like character string
\end{SeeAlso}
\begin{SeeAlso}{Util\_TableGet*()}
get a single (1-element array) value,
or more generally the first array element of an array value
\end{SeeAlso}
\begin{SeeAlso}{Util\_TableGet*Array()}
get an array value
\end{SeeAlso}
\begin{SeeAlso}{Util\_TableGetCharArray()}
get an array-of-\verb|CCTK_CHAR| value
\end{SeeAlso}
\begin{SeeAlso}{Util\_TableGetGeneric()}
get a single (1-element array) value with generic data type
\end{SeeAlso}
\begin{SeeAlso}{Util\_TableGetGenericArray()}
get an array value with generic data type
\end{SeeAlso}
\begin{SeeAlso}{Util\_TableSet*()}
set a single (1-element array) value
\end{SeeAlso}
\begin{SeeAlso}{Util\_TableSet*Array()}
set an array value
\end{SeeAlso}
\begin{SeeAlso}{Util\_TableSetGeneric()}
set a single (1-element array) value with generic data type
\end{SeeAlso}
\begin{SeeAlso}{Util\_TableSetGenericArray()}
set an array value with generic data type
\end{SeeAlso}
\begin{SeeAlso}{Util\_TableSetString()}
set a character-string value
\end{SeeAlso}
\begin{SeeAlso}{Util\_TableSetFromString()}
convenience routine to set key/value entries in a table based on a
parameter-file--like character string
\end{SeeAlso}
\begin{SeeAlso}{Util\_TableSetCharArray()}
set an array-of-\verb|CCTK_CHAR| value
\end{SeeAlso}
\end{SeeAlsoSection}

\begin{ErrorSection}
\begin{Error}{UTIL\_ERROR\_BAD\_HANDLE}
handle is invalid
\end{Error}
\begin{Error}{UTIL\_ERROR\_TABLE\_BAD\_KEY}
key contains '/' character
\end{Error}
\begin{Error}{UTIL\_ERROR\_BAD\_INPUT}
\verb|buffer != NULL| and \verb|buffer_length| $\le 0$
\end{Error}
\begin{Error}{UTIL\_ERROR\_TABLE\_NO\_SUCH\_KEY}
no such key in table
\end{Error}
\begin{Error}{UTIL\_ERROR\_TABLE\_WRONG\_DATA\_TYPE}
value has data type other than \verb|CCTK_CHAR|
\end{Error}
\begin{Error}{UTIL\_ERROR\_TABLE\_STRING\_TRUNCATED}
\quad
\verb|buffer != NULL| and value was truncated to fit in \verb|buffer[]|
\end{Error}
\end{ErrorSection}

\begin{ExampleSection}
\begin{Example}{C}
\begin{verbatim}
#include "util_ErrorCodes.h"
#include "util_Table.h"

#define N_BUFFER        100
char buffer[N_BUFFER];

int handle = Util_TableCreate(UTIL_TABLE_FLAGS_DEFAULT);
Util_TableSetString(handle, "relativity", "Einstein");

/* get length of string (= 10 here) */
int length = Util_TableGetString(handle, 0, NULL, "Einstein");

/* get null-terminated string into buffer, also returns 10 */
Util_TableGetString(handle, N_BUFFER, buffer, "Einstein");
\end{verbatim}
\end{Example}
\end{ExampleSection}
\end{FunctionDescription}

%%%%%%%%%%%%%%%%%%%%%%%%%%%%%%%%%%%%%%%%

\begin{FunctionDescription}{Util\_TableItAdvance}
\label{Util-TableItAdvance}
Advance a table iterator to the next entry in the table

\begin{SynopsisSection}
\begin{Synopsis}{C}
\begin{verbatim}
#include "util_ErrorCodes.h"
#include "util_Table.h"
int is_nonnull = Util_TableItAdvance(int ihandle);
\end{verbatim}
\end{Synopsis}
\end{SynopsisSection}

\begin{ResultSection}
\begin{Result}{\rm 1}
ok (iterator now points to some table entry)
\end{Result}
\begin{Result}{\rm 0}
ok (iterator has just advanced past the last table entry,
    and is now in the "null-pointer" state)
\end{Result}
\end{ResultSection}

\begin{ParameterSection}
\begin{Parameter}{ihandle ($\ge 0$)}
handle to the iterator
\end{Parameter}
\end{ParameterSection}

\begin{Discussion}
If we view an iterator as an abstraction of a pointer into the table,
then this function is the abstraction of the C \verb|++| operation applied
to the pointer, except that this function automagically sets the iterator
to the "null-pointer" state when it advances past the last table entry.

Note that bad things (garbage results, core dumps) may happen if
you call this function on an iterator which has been invalidated
by a change in the table's contents.
\end{Discussion}

\begin{ErrorSection}
\begin{Error}{UTIL\_ERROR\_BAD\_HANDLE}
iterator handle is invalid
\end{Error}
\end{ErrorSection}

\begin{ExampleSection}
\begin{Example}{C}
\begin{verbatim}
/* walk through all entries of a table */
int ihandle;

        for ( ihandle = Util_TableItCreate(handle) ;
              Util_TableItQueryIsNonNull(ihandle) > 0 ;
              Util_TableItAdvance(ihandle) )
        {
        /* do something with the table entry */
        }

Util_TableItDestroy(ihandle);
\end{verbatim}
\end{Example}
\end{ExampleSection}
\end{FunctionDescription}

%%%%%%%%%%%%%%%%%%%%%%%%%%%%%%%%%%%%%%%%

\begin{FunctionDescription}{Util\_TableItClone}
\label{Util-TableItClone}
Creates a new table iterator which is a ``clone'' (exact copy) of an
existing table iterator

\begin{SynopsisSection}
\begin{Synopsis}{C}
\begin{verbatim}
#include "util_ErrorCodes.h"
#include "util_Table.h"
int clone_ihandle = Util_TableItClone(int ihandle);
\end{verbatim}
\end{Synopsis}
\end{SynopsisSection}

\begin{ResultSection}
\begin{Result}{clone\_ihandle ($\ge 0$)}
A handle to the clone table iterator
\end{Result}
\end{ResultSection}

\begin{ParameterSection}
\begin{Parameter}{ihandle}
Handle to the table iterator to be cloned
\end{Parameter}
\end{ParameterSection}

\begin{Discussion}
This function creates a new iterator which points to the same place
in the same table as the original iterator.  If the original iterator
is in the ``null-pointer'' state, then the clone is also in this state.

Note that bad things (garbage results, core dumps) may happen if
you call this function on an iterator which has been invalidated
by a change in the table's contents.
\end{Discussion}

\begin{SeeAlsoSection}
\begin{SeeAlso}{Util\_TableClone()}
create a new table which is a ``clone'' (exact copy) of an existing
table
\end{SeeAlso}
\begin{SeeAlso}{Util\_TableItCreate()}
create a table iterator
\end{SeeAlso}
\begin{SeeAlso}{Util\_TableItDestroy()}
destroy a table iterator
\end{SeeAlso}
\end{SeeAlsoSection}

\begin{ErrorSection}
\begin{Error}{UTIL\_ERROR\_BAD\_HANDLE}
iterator handle to be cloned, is invalid
\end{Error}
\begin{Error}{UTIL\_ERROR\_NO\_MEMORY}
unable to allocate memory
\end{Error}
\end{ErrorSection}

\begin{ExampleSection}
\begin{Example}{C}
\begin{verbatim}
#include "util_ErrorCodes.h"
#include "util_Table.h"

/*
 * Apart from efficiency and slight differences in error return codes,
 * Util_TableItClone() could be simulated by the following code.
 */
int Util_TableItClone(int ihandle)
{
int status;

/* to what table does the to-be-cloned iterator point? */
const int handle = Util_TableQueryTableHandle(ihandle);
if (handle < 0)
        return handle;                /* error in querying table handle */

/* create the to-be-cloned iterator */
/* (pointing into the same table as the original iterator) */
  {
const int clone_ihandle = Util_TableItCreate(handle);
if (clone_ihandle < 0)
        return clone_ihandle;       /* error in creating clone iterator */

/* how long is the key to which the to-be-cloned iterator points? */
  {
const int key_length = Util_TableItQueryKeyValueInfo(ihandle,
                                                     0, NULL,
                                                     NULL, NULL);
if (key_length == UTIL_TABLE_ITERATOR_IS_NULL)
        {
        /* to-be-cloned iterator is in "null-pointer" state */
        Util_TableItSetToNull(clone_ihandle);
        return clone_ihandle;                          /* normal return */
        }
if (key_length < 0)
        return key_length;   /* error in querying to-be-cloned iterator */

/* to what key does the to-be-cloned iterator point? */
  {
const int key_buffer_length = key_length + 1;
char *const key_buffer = (char *) malloc(key_buffer_length);
if (key_buffer == NULL)
        return UTIL_ERROR_NO_MEMORY;
status = Util_TableItQueryKeyValueInfo(ihandle,
                                       key_buffer_length, key_buffer);
if (status < 0)
        return status;       /* error in querying to-be-cloned iterator */

/* set the clone iterator to point to the same key as the original */
status = Util_TableItSetToKey(clone_ihandle, key_buffer);
free(key_buffer);
return clone_ihandle;                                  /* normal return */
  }
  }
  }
}
\end{verbatim}
\end{Example}
\end{ExampleSection}
\end{FunctionDescription}

%%%%%%%%%%%%%%%%%%%%%%%%%%%%%%%%%%%%%%%%

\begin{FunctionDescription}{Util\_TableItCreate}
\label{Util-TableItCreate}
Create a new table iterator

\begin{SynopsisSection}
\begin{Synopsis}{C}
\begin{verbatim}
#include "util_ErrorCodes.h"
#include "util_Table.h"
int ihandle = Util_TableItCreate(int handle);
\end{verbatim}
\end{Synopsis}
\end{SynopsisSection}

\begin{ResultSection}
\begin{Result}{ihandle ($\ge 0$)}
handle to the iterator
\end{Result}
\end{ResultSection}

\begin{ParameterSection}
\begin{Parameter}{handle ($\ge 0$)}
handle to the table over which the iterator should iterate
\end{Parameter}
\end{ParameterSection}

\begin{Discussion}
This function creates a new table iterator.  The iterator initially
points at the starting table entry.
\end{Discussion}

\begin{SeeAlsoSection}
\begin{SeeAlso}{Util\_TableItDestroy()}
destroy a table iterator
\end{SeeAlso}
\end{SeeAlsoSection}

\begin{ErrorSection}
\begin{Error}{UTIL\_ERROR\_BAD\_HANDLE}
table handle is invalid
\end{Error}
\begin{Error}{UTIL\_ERROR\_NO\_MEMORY}
unable to allocate memory
\end{Error}
\end{ErrorSection}

\begin{ExampleSection}
\begin{Example}{C}
\begin{verbatim}
/* walk through all entries of a table */
int ihandle;

        for ( ihandle = Util_TableItCreate(handle) ;
              Util_TableItQueryIsNonNull(ihandle) > 0 ;
              Util_TableItAdvance(ihandle) )
        {
        /* do something with the table entry */
        }

Util_TableItDestroy(ihandle);
\end{verbatim}
\end{Example}
\end{ExampleSection}
\end{FunctionDescription}

%%%%%%%%%%%%%%%%%%%%%%%%%%%%%%%%%%%%%%%%

\begin{FunctionDescription}{Util\_TableItDestroy}
\label{Util-TableItDestroy}
Destroy a table iterator

\begin{SynopsisSection}
\begin{Synopsis}{C}
\begin{verbatim}
#include "util_ErrorCodes.h"
#include "util_Table.h"
int status = Util_TableItDestroy(int ihandle);
\end{verbatim}
\end{Synopsis}
\end{SynopsisSection}

\begin{ResultSection}
\begin{Result}{\rm 0}
ok
\end{Result}
\end{ResultSection}

\begin{ParameterSection}
\begin{Parameter}{ihandle ($\ge 0$)}
handle to the iterator
\end{Parameter}
\end{ParameterSection}

\begin{Discussion}
\end{Discussion}

\begin{SeeAlsoSection}
\begin{SeeAlso}{Util\_TableItCreate()}
create a table iterator
\end{SeeAlso}
\end{SeeAlsoSection}

\begin{ErrorSection}
\begin{Error}{UTIL\_ERROR\_BAD\_HANDLE}
iterator handle is invalid
\end{Error}
\begin{Error}{UTIL\_ERROR\_NO\_MEMORY}
unable to allocate memory
\end{Error}
\end{ErrorSection}

\begin{ExampleSection}
\begin{Example}{C}
\begin{verbatim}
/* walk through all entries of a table */
int ihandle;

        for ( ihandle = Util_TableItCreate(handle) ;
              Util_TableItQueryIsNonNull(ihandle) > 0 ;
              Util_TableItAdvance(ihandle) )
        {
        /* do something with the table entry */
        }

Util_TableItDestroy(ihandle);
\end{verbatim}
\end{Example}
\end{ExampleSection}
\end{FunctionDescription}

%%%%%%%%%%%%%%%%%%%%%%%%%%%%%%%%%%%%%%%%

\begin{FunctionDescription}{Util\_TableItQueryIsNonNull}
\label{Util-TableItQueryIsNonNull}
Query whether a table iterator is {\em not\/} in the ``null-pointer'' state

\begin{SynopsisSection}
\begin{Synopsis}{C}
\begin{verbatim}
#include "util_ErrorCodes.h"
#include "util_Table.h"
int status = Util_TableItQueryIsNonNull(int ihandle);
\end{verbatim}
\end{Synopsis}
\end{SynopsisSection}

\begin{ResultSection}
\begin{Result}{\rm 1}
iterator is {\em not\/} in the ``null-pointer'' state,
\ie{} iterator points to some table entry
\end{Result}
\begin{Result}{\rm 0}
iterator is in the ``null-pointer'' state
\end{Result}
\end{ResultSection}

\begin{ParameterSection}
\begin{Parameter}{ihandle ($\ge 0$)}
handle to the iterator
\end{Parameter}
\end{ParameterSection}

\begin{Discussion}
If no errors occur,
\verb|Util_TableItQueryIsNonNull(ihandle)|
is the same as\\
\verb|1 - Util_TableItQueryIsNull(ihandle)|.

Note that bad things (garbage results, core dumps) may happen if
you call this function on an iterator which has been invalidated
by a change in the table's contents.
\end{Discussion}

\begin{SeeAlsoSection}
\begin{SeeAlso}{Util\_TableItQueryIsNull()}
query whether a table iterator is in the ``null-pointer'' state
\end{SeeAlso}
\end{SeeAlsoSection}

\begin{ErrorSection}
\begin{Error}{UTIL\_ERROR\_BAD\_HANDLE}
iterator handle is invalid
\end{Error}
\end{ErrorSection}

\begin{ExampleSection}
\begin{Example}{C}
\begin{verbatim}
/* walk through all entries of a table */
int ihandle;

        for ( ihandle = Util_TableItCreate(handle) ;
              Util_TableItQueryIsNonNull(ihandle) > 0 ;
              Util_TableItAdvance(ihandle) )
        {
        /* do something with the table entry */
        }

Util_TableItDestroy(ihandle);
\end{verbatim}
\end{Example}
\end{ExampleSection}
\end{FunctionDescription}

%%%%%%%%%%%%%%%%%%%%%%%%%%%%%%%%%%%%%%%%

\begin{FunctionDescription}{Util\_TableItQueryIsNull}
\label{Util-TableItQueryIsNull}
Query whether a table iterator is in the ``null-pointer'' state

\begin{SynopsisSection}
\begin{Synopsis}{C}
\begin{verbatim}
#include "util_ErrorCodes.h"
#include "util_Table.h"
int status = Util_TableItQueryIsNull(int ihandle);
\end{verbatim}
\end{Synopsis}
\end{SynopsisSection}

\begin{ResultSection}
\begin{Result}{\rm 1}
iterator is in the ``null-pointer'' state
\end{Result}
\begin{Result}{\rm 0}
iterator is {\em not\/} in the ``null-pointer'' state,
\ie{} iterator points to some table entry
\end{Result}
\end{ResultSection}

\begin{ParameterSection}
\begin{Parameter}{ihandle ($\ge 0$)}
handle to the iterator
\end{Parameter}
\end{ParameterSection}

\begin{Discussion}
If no errors occur,
\verb|Util_TableItQueryIsNull(ihandle)|
is the same as
\verb|1 - Util_TableItQueryIsNonNull(ihandle)|.

Note that bad things (garbage results, core dumps) may happen if
you call this function on an iterator which has been invalidated
by a change in the table's contents.
\end{Discussion}

\begin{SeeAlsoSection}
\begin{SeeAlso}{Util\_TableItQueryIsNonNull()}
query whether a table iterator is {\em not\/} in the ``null-pointer'' state,
\ie{} whether the iterator points to some table entry
\end{SeeAlso}
\end{SeeAlsoSection}

\begin{ErrorSection}
\begin{Error}{UTIL\_ERROR\_BAD\_HANDLE}
iterator handle is invalid
\end{Error}
\end{ErrorSection}

\begin{ExampleSection}
\begin{Example}{C}
\begin{verbatim}
/* variant code to walk through all entries of a table */
int ihandle;

        for ( ihandle = Util_TableItCreate(handle) ;
              Util_TableItQueryIsNull(ihandle) == 0 ;
              Util_TableItAdvance(ihandle) )
        {
        /* do something with the table entry */
        }

Util_TableItDestroy(ihandle);
\end{verbatim}
\end{Example}
\end{ExampleSection}
\end{FunctionDescription}

%%%%%%%%%%%%%%%%%%%%%%%%%%%%%%%%%%%%%%%%

\begin{FunctionDescription}{Util\_TableItQueryKeyValueInfo}
\label{Util-TableItQueryKeyValueInfo}
Query the key and the type and number of elements of the value
corresponding to that key, of the table entry to which an iterator points

\begin{SynopsisSection}
\begin{Synopsis}{C}
\begin{verbatim}
#include "util_ErrorCodes.h"
#include "util_Table.h"
int key_length =
 Util_TableItQueryKeyValueInfo(int ihandle,
                               int key_buffer_length, char key_buffer[],
                               CCTK_INT *type_code, CCTK_INT *N_elements)
\end{verbatim}
\end{Synopsis}
\end{SynopsisSection}

\begin{ResultSection}
\begin{Result}{key\_length}
The string length of the key
(this has C {\t strlen()} semantics, \ie{} it does {\em not\/}
include a terminating null character)
\end{Result}
\end{ResultSection}

\begin{ParameterSection}
\begin{Parameter}{handle ($\ge 0$)}
handle to the table iterator
\end{Parameter}
\begin{Parameter}{key\_buffer\_length}
\quad
the length (\verb|sizeof()|) of \verb|key_buffer[]|
(must be $\ge 1$ if \verb|key_buffer != NULL|)
\end{Parameter}
\begin{Parameter}{key\_buffer}
a pointer to a buffer into which this function should store
(at most \verb|key_buffer_length-1| characters of) the key,
terminated by a null character as usual for C strings,
or NULL pointer to skip storing this
\end{Parameter}
\begin{Parameter}{type\_code}
a pointer to where this function should store the value's type code
(one of the \verb|CCTK_VARIABLE_|* constants from \verb|"cctk_Constants.h"|),
or a NULL pointer to skip storing this.
\end{Parameter}
\begin{Parameter}{N\_elements}
a pointer to where this function should store
the number of array elements in the value,
or a NULL pointer to skip storing this.
\end{Parameter}
\end{ParameterSection}

\begin{Discussion}
The usual use of an iterator is to iterate through all the entries
of a table, calling this function on each entry, then taking further
action based on the results.

Note that bad things (garbage results, core dumps) may happen if
you call this function on an iterator which has been invalidated
by a change in the table's contents.

If the error code \verb|UTIL_ERROR_TABLE_STRING_TRUNCATED| is returned,
then the first \verb|key_buffer_length-1| characters of the key are
returned in the user's key buffer (assuming \verb|key_buffer| is non-NULL),
followed by a null character to properly terminate the string in the
buffer.  If any other error code is returned, the user's key buffer
(pointed to by \verb|key_buffer| if this is non-NULL) is unchanged.
\end{Discussion}

\begin{SeeAlsoSection}
\begin{SeeAlso}{Util\_TableQueryValueInfo()}
query key present/absent in table, and optionally type and/or number
of elements, but using the key instead of an iterator
\end{SeeAlso}
\end{SeeAlsoSection}

\begin{ErrorSection}
\begin{Error}{UTIL\_ERROR\_BAD\_HANDLE}
handle is invalid
\end{Error}
\begin{Error}{UTIL\_ERROR\_TABLE\_ITERATOR\_IS\_NULL}
\quad
iterator is in "null-pointer" state
\end{Error}
\begin{Error}{UTIL\_ERROR\_TABLE\_STRING\_TRUNCATED}
\quad
\verb|key_buffer != NULL| and key was truncated to fit in \verb|key_buffer|
\end{Error}
\end{ErrorSection}

\begin{ExampleSection}
\begin{Example}{C}
\begin{verbatim}
/* print out all entries in a table */
/* return 0 for ok, type code for any types we can't handle, */
/*                  -ve for other errors */
#include <stdio.h>
#include <stdlib.h>
#include "util_ErrorCodes.h"
#include "util_Table.h"
#include "cctk.h"

int print_table(int handle)
{
int max_key_length, N_key_buffer, ihandle;
char *key_buffer;

max_key_length = Util_TableQueryMaxKeyLength(handle);
if (max_key_length < 0)
        return max_key_length;

N_key_buffer = max_key_length + 1;
key_buffer = (char *) malloc(N_key_buffer);
if (key_buffer == NULL)
        return UTIL_ERROR_NO_MEMORY;

        for ( ihandle = Util_TableItCreate(handle) ;
              Util_TableItQueryIsNonNull(ihandle) > 0 ;
              Util_TableItAdvance(ihandle) )
        {
        CCTK_INT type_code, N_elements;
        CCTK_INT value_int;
        CCTK_REAL value_real;

        Util_TableItQueryKeyValueInfo(ihandle,
                                      N_key_buffer, key_buffer,
                                      &type_code, &N_elements);
        printf("key = \"%s\"\n", key_buffer);

        switch  (type_code)
                {
        case CCTK_VARIABLE_INT:
                Util_TableGetInt(handle, &value_int, key_buffer);
                printf("value[int] = %d\n", (int)value_int);
                break;
        case CCTK_VARIABLE_REAL:
                Util_TableGetReal(handle, &value_real, key_buffer);
                printf("value[real] = %g\n", (double)value_real);
                break;
        default:
                /* we don't know how to handle this type */
                Util_TableItDestroy(ihandle);
                free(key_buffer);
                return type_code;
                }
        }

Util_TableItDestroy(ihandle);
free(key_buffer);
return 0;
}
\end{verbatim}
\end{Example}
\end{ExampleSection}
\end{FunctionDescription}

%%%%%%%%%%%%%%%%%%%%%%%%%%%%%%%%%%%%%%%%

\begin{FunctionDescription}{Util\_TableItQueryTableHandle}
\label{Util-TableItQueryTableHandle}
Query what table a table iterator iterates over

\begin{SynopsisSection}
\begin{Synopsis}{C}
\begin{verbatim}
#include "util_ErrorCodes.h"
#include "util_Table.h"
int handle = Util_TableItQueryTableHandle(int ihandle);
\end{verbatim}
\end{Synopsis}
\end{SynopsisSection}

\begin{ResultSection}
\begin{Result}{handle ($\ge 0$)}
handle to the table over which the iterator iterates
\end{Result}
\end{ResultSection}

\begin{ParameterSection}
\begin{Parameter}{ihandle ($\ge 0$)}
handle to the iterator
\end{Parameter}
\end{ParameterSection}

\begin{Discussion}
Note that it is always ok to call this function, regardless of
whether or not the iterator is in the ``null-pointer'' state.

It's also ok to call this function even when the iterator has been
invalidated by a change in the table's contents.
\end{Discussion}

\begin{SeeAlsoSection}
\begin{SeeAlso}{Util\_TableItCreate()}
create an iterator (which iterates over a specified table)
\end{SeeAlso}
\end{SeeAlsoSection}

\begin{ErrorSection}
\begin{Error}{UTIL\_ERROR\_BAD\_HANDLE}
iterator handle is invalid
\end{Error}
\end{ErrorSection}
\end{FunctionDescription}

%%%%%%%%%%%%%%%%%%%%%%%%%%%%%%%%%%%%%%%%

\begin{FunctionDescription}{Util\_TableItResetToStart}
\label{Util-TableItResetToStart}
Reset a table iterator to point to the starting table entry

\begin{SynopsisSection}
\begin{Synopsis}{C}
\begin{verbatim}
#include "util_ErrorCodes.h"
#include "util_Table.h"
int status = Util_TableItResetToStart(int ihandle);
\end{verbatim}
\end{Synopsis}
\end{SynopsisSection}

\begin{ResultSection}
\begin{ResultNote}
Results are the same as calling \verb|Util_TableItQueryIsNonNull()|
on the iterator after the reset:
\end{ResultNote}
\begin{Result}{\rm 1}
iterator is {\em not\/} in the ``null-pointer'' state,
\ie{} iterator points to some table entry
\end{Result}
\begin{Result}{\rm 0}
iterator is in the ``null-pointer'' state
(this happens if and only if the table is empty)
\end{Result}
\end{ResultSection}

\begin{ParameterSection}
\begin{Parameter}{ihandle ($\ge 0$)}
handle to the iterator
\end{Parameter}
\end{ParameterSection}

\begin{Discussion}
Note that it is always ok to call this function, regardless of
whether or not the iterator is in the ``null-pointer'' state.

It's also ok to call this function even when the iterator has been
invalidated by a change in the table's contents.
\end{Discussion}

\begin{SeeAlsoSection}
\begin{SeeAlso}{Util\_TableItSetToNull()}
set an iterator to the ``null-pointer'' state
\end{SeeAlso}
\begin{SeeAlso}{Util\_TableItSetToKey()}
set an iterator to point to a specified table entry
\end{SeeAlso}
\end{SeeAlsoSection}

\begin{ErrorSection}
\begin{Error}{UTIL\_ERROR\_BAD\_HANDLE}
iterator handle is invalid
\end{Error}
\end{ErrorSection}
\end{FunctionDescription}

%%%%%%%%%%%%%%%%%%%%%%%%%%%%%%%%%%%%%%%%

\begin{FunctionDescription}{Util\_TableItSetToKey}
\label{Util-TableItSetToKey}
Set a table iterator to point to a specified table entry

\begin{SynopsisSection}
\begin{Synopsis}{C}
\begin{verbatim}
#include "util_ErrorCodes.h"
#include "util_Table.h"
int status = Util_TableItSetToKey(int ihandle, const char *key);
\end{verbatim}
\end{Synopsis}
\end{SynopsisSection}

\begin{ResultSection}
\begin{Result}{\rm 0}
ok
\end{Result}
\end{ResultSection}

\begin{ParameterSection}
\begin{Parameter}{ihandle ($\ge 0$)}
handle to the iterator
\end{Parameter}
\end{ParameterSection}

\begin{Discussion}
This function has the same effect as calling \verb|Util_TableItResetToStart()|
followed by calling \verb|Util_TableItAdvance()| zero or more times
to make the iterator point to the desired table entry.  However, this
function will typically be (much) more efficient than that sequence.

Note that it is always ok to call this function, regardless of
whether or not the iterator is in the ``null-pointer'' state.

It's also ok to call this function even when the iterator has been
invalidated by a change in the table's contents.
\end{Discussion}

\begin{SeeAlsoSection}
\begin{SeeAlso}{Util\_TableItResetToStart()}
reset an iterator to point to the starting table entry
\end{SeeAlso}
\begin{SeeAlso}{Util\_TableItSetToNull()}
set a table iterator to the "null-pointer" state
\end{SeeAlso}
\end{SeeAlsoSection}

\begin{ErrorSection}
\begin{Error}{UTIL\_ERROR\_BAD\_HANDLE}
iterator handle is invalid
\end{Error}
\begin{Error}{UTIL\_ERROR\_TABLE\_BAD\_KEY}
key contains '/' character
\end{Error}
\begin{Error}{UTIL\_ERROR\_TABLE\_NO\_SUCH\_KEY}
no such key in table
\end{Error}
\end{ErrorSection}
\end{FunctionDescription}

%%%%%%%%%%%%%%%%%%%%%%%%%%%%%%%%%%%%%%%%

\begin{FunctionDescription}{Util\_TableItSetToNull}
\label{Util-TableItSetToNull}
Set a table iterator to the "null-pointer" state

\begin{SynopsisSection}
\begin{Synopsis}{C}
\begin{verbatim}
#include "util_ErrorCodes.h"
#include "util_Table.h"
int handle = Util_TableItSetToNull(int ihandle);
\end{verbatim}
\end{Synopsis}
\end{SynopsisSection}

\begin{ResultSection}
\begin{Result}{\rm 0}
ok
\end{Result}
\end{ResultSection}

\begin{ParameterSection}
\begin{Parameter}{ihandle ($\ge 0$)}
handle to the iterator
\end{Parameter}
\end{ParameterSection}

\begin{Discussion}
Note that it is always ok to call this function, regardless of
whether or not the iterator is already in the ``null-pointer'' state.

It's also ok to call this function even when the iterator has been
invalidated by a change in the table's contents.
\end{Discussion}

\begin{SeeAlsoSection}
\begin{SeeAlso}{Util\_TableItResetToStart()}
reset an iterator to point to the starting table entry
\end{SeeAlso}
\begin{SeeAlso}{Util\_TableItSetToKey()}
set an iterator to point to a specified table entry
\end{SeeAlso}
\end{SeeAlsoSection}

\begin{ErrorSection}
\begin{Error}{UTIL\_ERROR\_BAD\_HANDLE}
iterator handle is invalid
\end{Error}
\end{ErrorSection}
\end{FunctionDescription}

%%%%%%%%%%%%%%%%%%%%%%%%%%%%%%%%%%%%%%%%

\begin{FunctionDescription}{Util\_TableQueryFlags}
Query a table's flags word
\label{Util-TableQueryFlags}

\begin{SynopsisSection}
\begin{Synopsis}{C}
\begin{verbatim}
#include "util_ErrorCodes.h"
#include "util_Table.h"
int flags = Util_TableQueryFlags(int handle);
\end{verbatim}
\end{Synopsis}
\begin{Synopsis}{Fortran}
\begin{verbatim}
call Util_TableQueryFlags(flags, handle)
integer  flags, handle
\end{verbatim}
\end{Synopsis}
\end{SynopsisSection}

\begin{ResultSection}
\begin{Result}{flags ($\ge 0$)}
the flags word
\end{Result}
\end{ResultSection}

\begin{ParameterSection}
\begin{Parameter}{handle ($\ge 0$)}
handle to the table
\end{Parameter}
\end{ParameterSection}

\begin{Discussion}
See \verb|Util_TableCreate()| for further discussion of the semantics
of flag words.
\end{Discussion}

\begin{SeeAlsoSection}
\begin{SeeAlso}{Util\_TableClone()}
create a new table which is a ``clone'' (exact copy) of an existing
table
\end{SeeAlso}
\begin{SeeAlso}{Util\_TableCreate()}
create a table (flags word specified explicitly)
\end{SeeAlso}
\begin{SeeAlso}{Util\_TableCreateFromString()}
convenience routine to create a table (with certain default flags)
and set key/value entries in it based on a parameter-file--like
character string
\end{SeeAlso}
\end{SeeAlsoSection}

\begin{ErrorSection}
\begin{Error}{UTIL\_ERROR\_BAD\_HANDLE}
handle is invalid
\end{Error}
\end{ErrorSection}

\begin{ExampleSection}
\begin{Example}{C}
\begin{verbatim}
#include <string.h>
#include "util_ErrorCodes.h"
#include "util_String.h"
#include "util_Table.h"

/* compare two strings, doing the comparison with the same */
/* case-sensitive/insensitive semantics as a certain table uses */
int compare_strings(int handle, const char *str1, const char *str2)
{
int flags = Util_TableQueryFlags(handle);
return (flags & UTIL_TABLE_FLAGS_CASE_INSENSITIVE)
       ? Util_StrCmpi(str1, str2)
       :      strcmp (str1, str2);
}
\end{verbatim}
\end{Example}
\end{ExampleSection}
\end{FunctionDescription}

%%%%%%%%%%%%%%%%%%%%%%%%%%%%%%%%%%%%%%%%

\begin{FunctionDescription}{Util\_TableQueryValueInfo}
\label{Util-TableQueryValueInfo}
Query whether or not a specified key is in the table, and optionally
the type and/or number of elements of the value corresponding to this key

\begin{SynopsisSection}
\begin{Synopsis}{C}
\begin{verbatim}
#include "util_ErrorCodes.h"
#include "util_Table.h"
int key_exists =
 Util_TableQueryValueInfo(int handle,
                          CCTK_INT *type_code, CCTK_INT *N_elements,
                          const char *key);
\end{verbatim}
\end{Synopsis}
\begin{Synopsis}{Fortran}
\begin{verbatim}
call Util_TableQueryValueInfo(key_exists,
.                             handle,
.                             type_code, N_elements,
.                             key)
integer        key_exists, handle
CCTK_INT       type_code, N_elements
character*(*)  key
\end{verbatim}
\end{Synopsis}
\end{SynopsisSection}

\begin{ResultSection}
\begin{Result}{\rm 1}
ok (key is in table)
\end{Result}
\begin{Result}{\rm 0}
ok (no such key in table)\\
(in this case nothing is stored in \verb|*type_code| and \verb|*N_elements|)
\end{Result}
\end{ResultSection}

\begin{ParameterSection}
\begin{Parameter}{handle ($\ge 0$)}
handle to the table
\end{Parameter}
\begin{Parameter}{type\_code}
a pointer to where this function should store the value's type code
(one of the \verb|CCTK_VARIABLE_|* constants from \verb|"cctk_Constants.h"|),
or a NULL pointer to skip storing this.
\end{Parameter}
\begin{Parameter}{N\_elements}
a pointer to where this function should store
the number of array elements in the value,
or a NULL pointer to skip storing this.
\end{Parameter}
\begin{Parameter}{key}
a pointer to the key (a C-style null-terminated string)
\end{Parameter}
\end{ParameterSection}

\begin{Discussion}
Unlike all the other table query functions, this function returns
0 for ``no such key in table''.  The rationale for this design is that
by passing NULL pointers for \verb|type_code| and \verb|N_elements|, this
function is then a Boolean ``is key in table?'' predicate.

If any error code is returned, the user's buffers (pointed to by
\verb|type_code| and \verb|N_elements| if these are non-NULL) are
unchanged.
\end{Discussion}

\begin{SeeAlsoSection}
\begin{SeeAlso}{Util\_TableItQueryKeyValueInfo()}
query key present/absent in table, and optionally type and/or number
of elements, but using an iterator instead of the key
\end{SeeAlso}
\end{SeeAlsoSection}

\begin{ErrorSection}
\begin{Error}{UTIL\_ERROR\_BAD\_HANDLE}
handle is invalid
\end{Error}
\begin{Error}{UTIL\_ERROR\_TABLE\_BAD\_KEY}
key contains `/' character
\end{Error}
\end{ErrorSection}

\begin{ExampleSection}
\begin{Example}{C}
\begin{verbatim}
#include <stdio.h>
#include <assert.h>
#include "util_ErrorCodes.h"
#include "util_Table.h"

static const int data[] = {314, 159, 265};
#define N_DATA  (sizeof(data) / sizeof(data[0]))

CCTK_INT type_code, N_elements;

/* see whether or not "key" is in table */
if (Util_TableQueryValueInfo(handle, NULL, NULL, "key"))
        {
        /* key is in the table */
        }
   else {
        /* key is not in the table */
        }

/* put "data" in table as 3-element integer array */
Util_TableSetIntArray(handle, N_DATA, data, "data");

/* query info about "data" value */
assert( Util_TableQueryValueInfo(handle,
                                 &type_code, &N_elements,
                                 "data") == 1 );
assert( type_code == CCTK_VARIABLE_INT );
assert( N_elements == N_DATA );
\end{verbatim}
\end{Example}
\end{ExampleSection}
\end{FunctionDescription}

%%%%%%%%%%%%%%%%%%%%%%%%%%%%%%%%%%%%%%%%

\begin{FunctionDescription}{Util\_TableQueryMaxKeyLength}
\label{Util-TableQueryMaxKeyLength}
Query the maximum key length in a table

\begin{SynopsisSection}
\begin{Synopsis}{C}
\begin{verbatim}
#include "util_ErrorCodes.h"
#include "util_Table.h"
int max_key_length = Util_TableQueryMaxKeyLength(int handle);
\end{verbatim}
\end{Synopsis}
\begin{Synopsis}{Fortran}
\begin{verbatim}
call Util_TableQueryMaxKeyLength(max_key_length, handle)
integer  max_key_length, handle
\end{verbatim}
\end{Synopsis}
\end{SynopsisSection}

\begin{ResultSection}
\begin{Result}{max\_key\_length ($\ge 0$)}
The string length of the longest key in the table
(this has C {\t strlen()} semantics, \ie{} it does {\em not\/}
include a terminating null character)
\end{Result}
\end{ResultSection}

\begin{ParameterSection}
\begin{Parameter}{handle ($\ge 0$)}
handle to the table
\end{Parameter}
\end{ParameterSection}

\begin{Discussion}
This function is useful if you're going to iterate through a table,
and you want to allocate a buffer which is guaranteed to be big enough
to hold any key in the table.
\end{Discussion}

\begin{ErrorSection}
\begin{Error}{UTIL\_ERROR\_BAD\_HANDLE}
handle is invalid
\end{Error}
\end{ErrorSection}

\begin{ExampleSection}
\begin{Example}{C}
\begin{verbatim}
#include <stdlib.h>
#include "util_ErrorCodes.h"
#include "util_Table.h"
#include "cctk.h"

int max_key_length = Util_TableQueryMaxKeyLength(handle);
int N_buffer = max_key_length + 1;
char *const buffer = (char *) malloc(N_buffer);
if (buffer == NULL)
        {
        CCTK_WARN(-1, "couldn't allocate memory for table key buffer!");
        abort();        /* CCTK_Abort() would be better */
                        /* if we have a cGH* available */
        }

/* now buffer is guaranteed to be */
/* big enough for any key in the table */
\end{verbatim}
\end{Example}
\end{ExampleSection}
\end{FunctionDescription}

%%%%%%%%%%%%%%%%%%%%%%%%%%%%%%%%%%%%%%%%

\begin{FunctionDescription}{Util\_TableQueryNKeys}
\label{Util-TableQueryNKeys}
Query the number of key/value entries in a table

\begin{SynopsisSection}
\begin{Synopsis}{C}
\begin{verbatim}
#include "util_ErrorCodes.h"
#include "util_Table.h"
int N_Keys = Util_TableQueryNKeys(int handle);
\end{verbatim}
\end{Synopsis}
\begin{Synopsis}{Fortran}
\begin{verbatim}
call Util_TableQueryNKeys(N_Keys, handle)
integer  N_Keys, handle
\end{verbatim}
\end{Synopsis}
\end{SynopsisSection}

\begin{ResultSection}
\begin{Result}{N\_Keys ($\ge 0$)}
the number of key/value entries in the table
\end{Result}
\end{ResultSection}

\begin{ParameterSection}
\begin{Parameter}{handle ($\ge 0$)}
handle to the table
\end{Parameter}
\end{ParameterSection}

\begin{ErrorSection}
\begin{Error}{UTIL\_ERROR\_BAD\_HANDLE}
handle is invalid
\end{Error}
\end{ErrorSection}
\end{FunctionDescription}

%%%%%%%%%%%%%%%%%%%%%%%%%%%%%%%%%%%%%%%%

\begin{FunctionDescription}{Util\_TableSet*}
\label{Util-TableSet*}
This is a family of functions, one for each Cactus data type,
to set the value associated with a specified key to be a specified
single (1-element array) value

\begin{SynopsisSection}
\begin{Synopsis}{C}
\begin{verbatim}
#include "util_ErrorCodes.h"
#include "util_Table.h"
int status = Util_TableSetXxx(int handle,
                              CCTK_XXX value,
                              const char *key);
\end{verbatim}
where \verb|XXX| is one of
   \verb|POINTER|, \verb|FPOINTER|%%%
\footnote{%%%
	 For backwards compatability the function
	 {\tt Util\_TableSetFnPointer()} is also provided
	 as an alias for {\tt Util\_TableSetFPointer()}.
	 This is deprecated as of Cactus 4.0 beta 13.
	 }%%%
,
   \verb|CHAR|,
   \verb|INT|, \verb|INT2|, \verb|INT4|, \verb|INT8|,
   \verb|REAL|, \verb|REAL4|, \verb|REAL8|, \verb|REAL16|,
   \verb|COMPLEX|, \verb|COMPLEX8|, \verb|COMPLEX16|, \verb|COMPLEX32|
(not all of these may be supported on any given system)
\end{Synopsis}
\begin{Synopsis}{Fortran}
\begin{verbatim}
call Util_TableSetXxx(status, handle, value, key)
integer        status, handle
CCTK_XXX       value
character*(*)  key
\end{verbatim}
\end{Synopsis}
\end{SynopsisSection}

\begin{ResultSection}
\begin{Result}{\rm 1}
ok (key was already in table before this call, old value was replaced)\\
   (it doesn't matter what the old value's \verb|type_code| and
    \verb|N_elements| were, \ie{} these do {\em not\/} have to match
    the new value)
\end{Result}
\begin{Result}{\rm 0}
ok (key was not in table before this call)
\end{Result}
\end{ResultSection}

\begin{ParameterSection}
\begin{Parameter}{handle ($\ge 0$)}
handle to the table
\end{Parameter}
\begin{Parameter}{value}
the value to be associated with the key
\end{Parameter}
\begin{Parameter}{key}
a pointer to the key (a C-style null-terminated string)
\end{Parameter}
\end{ParameterSection}

\begin{Discussion}
The key may be any C character string which does not contain a slash
character (\verb|'/'|).

The value is stored as a 1-element array.

This function invalidates any iterators for the table which are
not in the ``null-pointer'' state.
\end{Discussion}

\begin{SeeAlsoSection}
\begin{SeeAlso}{Util\_TableCreateFromString()}
convenience routine to create a table and set key/value entries
in it based on a parameter-file--like character string
\end{SeeAlso}
\begin{SeeAlso}{Util\_TableGet*()}
get a single (1-element array) value,
or more generally the first array element of an array value
\end{SeeAlso}
\begin{SeeAlso}{Util\_TableGet*Array()}
get an array value
\end{SeeAlso}
\begin{SeeAlso}{Util\_TableGetGeneric()}
get a single (1-element array) value with generic data type
\end{SeeAlso}
\begin{SeeAlso}{Util\_TableGetGenericArray()}
get an array value with generic data type
\end{SeeAlso}
\begin{SeeAlso}{Util\_TableGet*String()}
get a character-string value
\end{SeeAlso}
\begin{SeeAlso}{Util\_TableSet*Array()}
set an array value
\end{SeeAlso}
\begin{SeeAlso}{Util\_TableSetGeneric()}
set a single (1-element array) value with generic data type
\end{SeeAlso}
\begin{SeeAlso}{Util\_TableSetGenericArray()}
set an array value with generic data type
\end{SeeAlso}
\begin{SeeAlso}{Util\_TableSetFromString()}
convenience routine to set key/value entries in a table based on a
parameter-file--like character string
\end{SeeAlso}
\begin{SeeAlso}{Util\_TableSetString()}
set a character-string value
\end{SeeAlso}
\end{SeeAlsoSection}

\begin{ErrorSection}
\begin{Error}{UTIL\_ERROR\_BAD\_HANDLE}
handle is invalid
\end{Error}
\begin{Error}{UTIL\_ERROR\_TABLE\_BAD\_KEY}
key contains '/' character
\end{Error}
\begin{Error}{UTIL\_ERROR\_NO\_MEMORY}
unable to allocate memory
\end{Error}
\end{ErrorSection}

\begin{ExampleSection}
\begin{Example}{C}
\begin{verbatim}
#include <math.h>
#include "util_ErrorCodes.h"
#include "util_Table.h"

CCTK_COMPLEX16 z;
int handle = Util_TableCreate(UTIL_TABLE_FLAGS_DEFAULT);

Util_TableSetInt(handle, 42, "the answer");
Util_TableSetReal(handle, 299792458.0, "speed of light");

z.Re = cos(0.37);       z.Im = sin(0.37);
Util_TableSetComplex16(handle, z, "my complex number");
\end{verbatim}
\end{Example}
\end{ExampleSection}
\end{FunctionDescription}

%%%%%%%%%%%%%%%%%%%%%%%%%%%%%%%%%%%%%%%%

\begin{FunctionDescription}{Util\_TableSet*Array}
\label{Util-TableSet*Array}
This is a family of functions, one for each Cactus data type,
to set the value associated with a specified key to be a copy
of a specified array

\begin{SynopsisSection}
\begin{Synopsis}{C}
\begin{verbatim}
#include "util_ErrorCodes.h"
#include "util_Table.h"
int status = Util_TableSetXxxArray(int handle,
                                   int N_elements,
                                   const CCTK_XXX array[],
                                   const char *key);
\end{verbatim}
where \verb|XXX| is one of
   \verb|POINTER|, \verb|FPOINTER|%%%
\footnote{%%%
	 For backwards compatability the function
	 {\tt Util\_TableSetFnPointerArray()} is also provided
	 as an alias for {\tt Util\_TableSetFPointerArray()}.
	 This is deprecated as of Cactus 4.0 beta 13.
	 }%%%
,
   \verb|CHAR|,
   \verb|INT|, \verb|INT2|, \verb|INT4|, \verb|INT8|,
   \verb|REAL|, \verb|REAL4|, \verb|REAL8|, \verb|REAL16|,
   \verb|COMPLEX|, \verb|COMPLEX8|, \verb|COMPLEX16|, \verb|COMPLEX32|
(not all of these may be supported on any given system)
\end{Synopsis}
\begin{Synopsis}{Fortran}
\begin{verbatim}
call Util_TableSetXxxArray(status, handle, N_elements, array, key)
integer        status, handle, N_elements
CCTK_XXX(*)    array
character*(*)  key
\end{verbatim}
\end{Synopsis}
\end{SynopsisSection}

\begin{ResultSection}
\begin{Result}{\rm 1}
ok (key was already in table before this call, old value was replaced)\\
   (it doesn't matter what the old value's \verb|type_code| and
    \verb|N_elements| were, \ie{} these do {\em not\/} have to match
    the new value)
\end{Result}
\begin{Result}{\rm 0}
ok (key was not in table before this call)
\end{Result}
\end{ResultSection}

\begin{ParameterSection}
\begin{Parameter}{handle ($\ge 0$)}
handle to the table
\end{Parameter}
\begin{Parameter}{N\_elements ($\ge 0$)}
the number of array elements in \verb|array[]|
\end{Parameter}
\begin{Parameter}{array}
a pointer to the array (a copy of which) is to be associated with the key
\end{Parameter}
\begin{Parameter}{key}
a pointer to the key (a C-style null-terminated string)
\end{Parameter}
\end{ParameterSection}

\begin{Discussion}
The key may be any C character string which does not contain a slash
character (\verb|'/'|).

Note that empty (0-element) arrays are ok.

This function invalidates any iterators for the table which are
not in the ``null-pointer'' state.
\end{Discussion}

\begin{SeeAlsoSection}
\begin{SeeAlso}{Util\_TableCreateFromString()}
convenience routine to create a table and set key/value entries
in it based on a parameter-file--like character string
\end{SeeAlso}
\begin{SeeAlso}{Util\_TableGet*()}
get a single (1-element array) value,
or more generally the first array element of an array value
\end{SeeAlso}
\begin{SeeAlso}{Util\_TableGet*Array()}
get an array value
\end{SeeAlso}
\begin{SeeAlso}{Util\_TableGetGeneric()}
get a single (1-element array) value with generic data type
\end{SeeAlso}
\begin{SeeAlso}{Util\_TableGetGenericArray()}
get an array value with generic data type
\end{SeeAlso}
\begin{SeeAlso}{Util\_TableGet*String()}
get a character-string value
\end{SeeAlso}
\begin{SeeAlso}{Util\_TableSet*()}
set a single (1-element array) value
\end{SeeAlso}
\begin{SeeAlso}{Util\_TableSetGeneric()}
set a single (1-element array) value with generic data type
\end{SeeAlso}
\begin{SeeAlso}{Util\_TableSetGenericArray()}
set an array value with generic data type
\end{SeeAlso}
\begin{SeeAlso}{Util\_TableSetFromString()}
convenience routine to set key/value entries in a table based on a
parameter-file--like character string
\end{SeeAlso}
\begin{SeeAlso}{Util\_TableSetString()}
set a character-string value
\end{SeeAlso}
\end{SeeAlsoSection}

\begin{ErrorSection}
\begin{Error}{UTIL\_ERROR\_BAD\_HANDLE}
handle is invalid
\end{Error}
\begin{Error}{UTIL\_ERROR\_TABLE\_BAD\_KEY}
key contains '/' character
\end{Error}
\begin{Error}{UTIL\_ERROR\_BAD\_INPUT}
\verb|N_elements| $< 0$
\end{Error}
\begin{Error}{UTIL\_ERROR\_NO\_MEMORY}
unable to allocate memory
\end{Error}
\end{ErrorSection}

\begin{ExampleSection}
\begin{Example}{C}
\begin{verbatim}
#include "util_ErrorCodes.h"
#include "util_Table.h"

#define N_DIGITS        5
static const CCTK_INT pi_digits[N_DIGITS] = {3, 14, 159, 2653, 58979};
int handle = Util_TableCreate(UTIL_TABLE_FLAGS_DEFAULT);

Util_TableSetIntArray(handle, N_DIGITS, pi_digits, "digits of pi");
\end{verbatim}
\end{Example}
\end{ExampleSection}
\end{FunctionDescription}

%%%%%%%%%%%%%%%%%%%%%%%%%%%%%%%%%%%%%%%%

\begin{FunctionDescription}{Util\_TableSetFromString}
\label{Util-TableSetFromString}
Sets values in a table based on a string argument, which is interpreted
with ``parameter-file'' semantics

\begin{SynopsisSection}
\begin{Synopsis}{C}
\begin{verbatim}
#include "util_ErrorCodes.h"
#include "util_Table.h"
int count = Util_TableSetFromString(int handle, const char *string);
\end{verbatim}
\end{Synopsis}
\begin{Synopsis}{Fortran}
\begin{verbatim}
call Util_TableSetFromString(count, handle, string)
integer        count, handle
character*(*)  string
\end{verbatim}
\end{Synopsis}
\end{SynopsisSection}

\begin{ResultSection}
\begin{Result}{count ($\ge 0$)}
the number of key/value entries set
\end{Result}
\end{ResultSection}

\begin{ParameterSection}
\begin{Parameter}{string}
a pointer to a C-style null-terminated string specifying the table
entries to be set (see below for details on the string contents)
\end{Parameter}
\end{ParameterSection}

\begin{Discussion}
The string should contain a sequence of zero or more \verb|key=value|
``assignments'', separated by whitespace.  This function processes
these assignments in left-to-right order, setting corresponding key/value
entries in the table.

The present implementation only recognises integer, real, and character-string
values (not complex), and integer and real arrays.
To be precise, the string must match the following BNF:\\
\quad
\begin{tabular}{l@{\quad$\rightarrow$\quad}p{9cm}}
string  & assign*                                                       \\
assign  & whitespace*                                                   \\
assign  & whitespace* key whitespace* = whitespace* value delimiter     \\
key     & any string not containing \verb|'/'| or \verb|'='| or whitespace\\
value   & array $\big|$ int\_value $\big|$ real\_value $\big|$ string\_value\\
array   & \{ int\_value* \} $\big|$ \{ real\_value \}                   \\
int\_value      & anything recognized as a valid integer by strdol(3)
                  in base 10                                            \\
real\_value     & anything not recognized as a valid integer by strtol(3)
                  but recognized as valid by strdod(3)                  \\
string\_value   & a C-style string enclosed in "double quotes"
                  (C-style character escape codes are allowed
                  ie. bell~(\verb|'\a'|), backspace~(\verb|'\b'|),
                  form-feed~(\verb|'\f'|), newline~(\verb|'\n'|),
                  carriage-return~(\verb|'\r'|), tab~(\verb|'\t'|),
                  vertical-tab~(\verb|'\v'|), backslash~(\verb|'\\'|),
                  single-quote~(\verb|'\''|), double-quote~(\verb|'\"'|),
                  question-mark~(\verb|'\?'|))                          \\
string\_value   & a C-style string enclosed in 'single quotes'
                  (C-style character escape codes are {\bf not} allowed,
                  ie. every character within the string is interpreted
                  literally)                                            \\
delimiter       & end-of-string $\big|$ whitespace                      \\
whitespace      & blank~(\verb|' '|)
                  $\big|$ tab~(\verb|'\t'|)
                  $\big|$ newline~(\verb|'\n'|)
                  $\big|$ carriage-return~(\verb|'\r'|)
                  $\big|$ form-feed~(\verb|'\f'|)
                  $\big|$ vertical-tab~(\verb|'\v'|)                    %%%\\
\end{tabular}\\
where * denotes 0 or more repetitions and $\big|$ denotes logical or.

Notice also that the keys allowed by this function are somewhat more
restricted than those allowed by the other \verb|Util_TableSet*()|
functions, in that this function disallows keys containing \verb|'='|
and/or whitespace.

If any error code is returned, assignments lexicographically earlier
in the input string than where the error was detected will already
have been made in the table.  Unfortunately, there is no easy way to
find out where the error was detected. :(
\end{Discussion}

\begin{SeeAlsoSection}
\begin{SeeAlso}{Util\_TableCreateFromString()}
convenience routine to create a table and set key/value entries
in it based on a parameter-file--like character string
\end{SeeAlso}
\begin{SeeAlso}{Util\_TableGet*()}
get a single (1-element array) value,
or more generally the first array element of an array value
\end{SeeAlso}
\begin{SeeAlso}{Util\_TableGet*Array()}
get an array value
\end{SeeAlso}
\begin{SeeAlso}{Util\_TableGetGeneric()}
get a single (1-element array) value with generic data type
\end{SeeAlso}
\begin{SeeAlso}{Util\_TableGetGenericArray()}
get an array value with generic data type
\end{SeeAlso}
\begin{SeeAlso}{Util\_TableGet*String()}
get a character-string value
\end{SeeAlso}
\begin{SeeAlso}{Util\_TableSet*()}
set a single (1-element array) value
\end{SeeAlso}
\begin{SeeAlso}{Util\_TableSet*Array()}
set an array value
\end{SeeAlso}
\begin{SeeAlso}{Util\_TableSetGeneric()}
set a single (1-element array) value with generic data type
\end{SeeAlso}
\begin{SeeAlso}{Util\_TableSetGenericArray()}
set an array value with generic data type
\end{SeeAlso}
\begin{SeeAlso}{Util\_TableSetString()}
set a character-string value
\end{SeeAlso}
\end{SeeAlsoSection}

\begin{ErrorSection}
\begin{Error}{UTIL\_ERROR\_NO\_MEMORY}
unable to allocate memory
\end{Error}
\begin{Error}{UTIL\_ERROR\_BAD\_KEY}
invalid input: key contains invalid character
\end{Error}
\begin{Error}{UTIL\_ERROR\_BAD\_INPUT}
invalid input: can't parse input string
\end{Error}
\begin{Error}{UTIL\_ERROR\_NO\_MIXED\_TYPE\_ARRAY}
invalid input: different array values have different datatypes
\end{Error}
\begin{Error}{\rm other error codes}
this function may also return any error codes returned by
\verb|Util_TableSetString()|, \verb|Util_TableSetInt()|,
\verb|Util_TableSetReal()|, \verb|Util_TableSetIntArray()|,
or \verb|Util_TableSetRealArray()|.
\end{Error}
\end{ErrorSection}

\begin{ExampleSection}
\begin{Example}{C}
\begin{verbatim}
#include "util_ErrorCodes.h"
#include "util_Table.h"

/* suppose we have a table referred to by  handle */

/* then the call... */
int count = Util_TableSetFromString(handle, "n = 6\t"
                                            "dx = 4.0e-5\t"
                                            "pi = 3.1\t"
                                            "s = 'my string'\t"
                                            "array = { 1 2 3 }");
/* ... will return count=5 ... */

/* ... and is otherwise equivalent to the five calls ... */
CCTK_INT array[] = {1, 2, 3};

Util_TableSetInt(handle, 6, "n");
Util_TableSetReal(handle, 4.0e-5, "dx");
Util_TableSetReal(handle, 3.1, "pi");
Util_TableSetString(handle, "my string", "s");
Util_TableSetIntArray(handle, 3, array, "array");
\end{verbatim}
\end{Example}
\end{ExampleSection}
\end{FunctionDescription}

%%%%%%%%%%%%%%%%%%%%%%%%%%%%%%%%%%%%%%%%

\begin{FunctionDescription}{Util\_TableSetGeneric}
\label{Util-TableSetGeneric}
Set the value associated with a specified key to be a specified
single (1-element array) value, whose data type is generic.  That
is, the value is specified by a \verb|CCTK_VARIABLE_|* type code
and a \verb|void *| pointer.

\begin{SynopsisSection}
\begin{Synopsis}{C}
\begin{verbatim}
#include "util_ErrorCodes.h"
#include "util_Table.h"
int status = Util_TableSetGeneric(int handle,
                                  int type_code, const void *value,
                                  const char *key);
\end{verbatim}
\end{Synopsis}
\begin{Synopsis}{Fortran}
\begin{verbatim}
call Util_TableSetGeneric(status, handle, type_code, value, key)
integer        status, handle, type_code
CCTK_POINTER   value
character*(*)  key
\end{verbatim}
\end{Synopsis}
\end{SynopsisSection}

\begin{ResultSection}
\begin{Result}{\rm 1}
ok (key was already in table before this call, old value was replaced)\\
   (it doesn't matter what the old value's \verb|type_code| and
    \verb|N_elements| were, \ie{} these do {\em not\/} have to match
    the new value)
\end{Result}
\begin{Result}{\rm 0}
ok (key was not in table before this call)
\end{Result}
\end{ResultSection}

\begin{ParameterSection}
\begin{Parameter}{handle ($\ge 0$)}
handle to the table
\end{Parameter}
\begin{Parameter}{type\_code}
the array elements' type code
(one of the \verb|CCTK_VARIABLE_|* constants from \verb|"cctk_Constants.h"|)
\end{Parameter}
\begin{Parameter}{value\_ptr}
a pointer to the value to be associated with the key
\end{Parameter}
\begin{Parameter}{key}
a pointer to the key (a C-style null-terminated string)
\end{Parameter}
\end{ParameterSection}

\begin{Discussion}
The key may be any C character string which does not contain a slash
character (\verb|'/'|).

The value is stored as a 1-element array.

This function invalidates any iterators for the table which are
not in the ``null-pointer'' state.
\end{Discussion}

\begin{SeeAlsoSection}
\begin{SeeAlso}{Util\_TableCreateFromString()}
convenience routine to create a table and set key/value entries
in it based on a parameter-file--like character string
\end{SeeAlso}
\begin{SeeAlso}{Util\_TableGet*()}
get a single (1-element array) value,
or more generally the first array element of an array value
\end{SeeAlso}
\begin{SeeAlso}{Util\_TableGet*Array()}
get an array value
\end{SeeAlso}
\begin{SeeAlso}{Util\_TableGetGeneric()}
get a single (1-element array) value with generic data type
\end{SeeAlso}
\begin{SeeAlso}{Util\_TableGetGenericArray()}
get an array value with generic data type
\end{SeeAlso}
\begin{SeeAlso}{Util\_TableGet*String()}
get a character-string value
\end{SeeAlso}
\begin{SeeAlso}{Util\_TableSet*()}
set a single (1-element array) value
\end{SeeAlso}
\begin{SeeAlso}{Util\_TableSet*Array()}
set an array value
\end{SeeAlso}
\begin{SeeAlso}{Util\_TableSetGeneric()}
set a single (1-element array) value with generic data type
\end{SeeAlso}
\begin{SeeAlso}{Util\_TableSetGenericArray()}
set an array value with generic data type
\end{SeeAlso}
\begin{SeeAlso}{Util\_TableSetFromString()}
convenience routine to set key/value entries in a table based on a
parameter-file--like character string
\end{SeeAlso}
\begin{SeeAlso}{Util\_TableSetString()}
set a character-string value
\end{SeeAlso}
\end{SeeAlsoSection}

\begin{ErrorSection}
\begin{Error}{UTIL\_ERROR\_BAD\_HANDLE}
handle is invalid
\end{Error}
\begin{Error}{UTIL\_ERROR\_TABLE\_BAD\_KEY}
key contains '/' character
\end{Error}
\begin{Error}{UTIL\_ERROR\_NO\_MEMORY}
unable to allocate memory
\end{Error}
\end{ErrorSection}

\begin{ExampleSection}
\begin{Example}{C}
\begin{verbatim}
#include "util_Table.h"
#include "cctk_Constants.h"

const CCTK_INT i = 42;
const void *iptr = (void *) &i;
CCTK_INT icopy;

const CCTK_REAL x = 299792458.0;
const void *xptr = (void *) &x;
CCTK_REAL xcopy;

const int handle = Util_TableCreate(UTIL_TABLE_FLAGS_DEFAULT);

Util_TableSetGeneric(handle, CCTK_VARIABLE_INT, iptr, "the answer");
Util_TableSetGeneric(handle, CCTK_VARIABLE_REAL, xptr, "speed of light");

/* gets icopy to 42 */
Util_TableGetInt(handle, &icopy, "the answer");

/* gets xcopy to 299792458.0 */
Util_TableGetReal(handle, &xcopy, "speed of light");
\end{verbatim}
\end{Example}
\end{ExampleSection}
\end{FunctionDescription}

%%%%%%%%%%%%%%%%%%%%%%%%%%%%%%%%%%%%%%%%

\begin{FunctionDescription}{Util\_TableSetGenericArray}
\label{Util-TableSetGenericArray}
Set the value associated with a specified key to be a copy of a
specified array, whose data type is generic.  That is, the array
is specified by a \verb|CCTK_VARIABLE_|* type code, a count of
the number of array elements, and a \verb|void *| pointer.

\begin{SynopsisSection}
\begin{Synopsis}{C}
\begin{verbatim}
#include "util_ErrorCodes.h"
#include "util_Table.h"
int status = Util_TableSetGenericArray(int handle,
                                       int type_code,
                                       int N_elements, const void *array,
                                       const char *key);
\end{verbatim}
\end{Synopsis}
\begin{Synopsis}{Fortran}
\begin{verbatim}
call Util_TableSetGenericArray(status,
.                              handle,
.                              type_code,
.                              N_elements, array,
.                              key)
integer           status, handle, type_code, N_elements
CCTK_POINTER(*)   array
character*(*)     key
\end{verbatim}
\end{Synopsis}
\end{SynopsisSection}

\begin{ResultSection}
\begin{Result}{\rm 1}
ok (key was already in table before this call, old value was replaced)\\
   (it doesn't matter what the old value's \verb|type_code| and
    \verb|N_elements| were, \ie{} these do {\em not\/} have to match
    the new value)
\end{Result}
\begin{Result}{\rm 0}
ok (key was not in table before this call)
\end{Result}
\end{ResultSection}

\begin{ParameterSection}
\begin{Parameter}{handle ($\ge 0$)}
handle to the table
\end{Parameter}
\begin{Parameter}{type\_code}
the array elements' type code
(one of the \verb|CCTK_VARIABLE_|* constants from \verb|"cctk_Constants.h"|)
\end{Parameter}
\begin{Parameter}{N\_elements ($\ge 0$)}
the number of array elements in \verb|array[]|
\end{Parameter}
\begin{Parameter}{value\_ptr}
a pointer to the value to be associated with the key
\end{Parameter}
\begin{Parameter}{key}
a pointer to the key (a C-style null-terminated string)
\end{Parameter}
\end{ParameterSection}

\begin{Discussion}
The key may be any C character string which does not contain a slash
character (\verb|'/'|).

The value is stored as a 1-element array.

This function invalidates any iterators for the table which are
not in the ``null-pointer'' state.
\end{Discussion}

\begin{SeeAlsoSection}
\begin{SeeAlso}{Util\_TableCreateFromString()}
convenience routine to create a table and set key/value entries
in it based on a parameter-file--like character string
\end{SeeAlso}
\begin{SeeAlso}{Util\_TableGet*()}
get a single (1-element array) value,
or more generally the first array element of an array value
\end{SeeAlso}
\begin{SeeAlso}{Util\_TableGet*Array()}
get an array value
\end{SeeAlso}
\begin{SeeAlso}{Util\_TableGetGeneric()}
get a single (1-element array) value with generic data type
\end{SeeAlso}
\begin{SeeAlso}{Util\_TableGetGenericArray()}
get an array value with generic data type
\end{SeeAlso}
\begin{SeeAlso}{Util\_TableGet*String()}
get a character-string value
\end{SeeAlso}
\begin{SeeAlso}{Util\_TableSet*()}
set a single (1-element array) value
\end{SeeAlso}
\begin{SeeAlso}{Util\_TableSet*Array()}
set an array value
\end{SeeAlso}
\begin{SeeAlso}{Util\_TableSetGeneric()}
set a single (1-element array) value with generic data type
\end{SeeAlso}
\begin{SeeAlso}{Util\_TableSetFromString()}
convenience routine to set key/value entries in a table based on a
parameter-file--like character string
\end{SeeAlso}
\begin{SeeAlso}{Util\_TableSetString()}
set a character-string value
\end{SeeAlso}
\end{SeeAlsoSection}

\begin{ErrorSection}
\begin{Error}{UTIL\_ERROR\_BAD\_HANDLE}
handle is invalid
\end{Error}
\begin{Error}{UTIL\_ERROR\_TABLE\_BAD\_KEY}
key contains '/' character
\end{Error}
\begin{Error}{UTIL\_ERROR\_NO\_MEMORY}
unable to allocate memory
\end{Error}
\end{ErrorSection}

\begin{ExampleSection}
\begin{Example}{C}
\begin{verbatim}
#include "util_Table.h"
#include "cctk_Constants.h"

#define N_IARRAY    3
const CCTK_INT iarray[N_IARRAY] = {42, 69, 105};
const void *iarray_ptr = (void *) iarray;
CCTK_INT iarray2[N_IARRAY];

#define N_XARRAY    2
const CCTK_REAL xarray[N_XARRAY] = {6.67e-11, 299792458.0};
const void *xarray_ptr = (void *) xarray;
CCTK_REAL xarray2[N_XARRAY];

const int handle = Util_TableCreate(UTIL_TABLE_FLAGS_DEFAULT);

Util_TableSetGenericArray(handle,
                          CCTK_VARIABLE_INT,
                          N_IARRAY, iarray_ptr,
                          "my integer array");
Util_TableSetGenericArray(handle,
                          CCTK_VARIABLE_REAL,
                          N_XARRAY, xarray_ptr,
                          "my real array");

/* gets iarray2[0] = 42, iarray2[1] = 69, iarray2[2] = 105 */
Util_TableGetIntArray(handle, N_IARRAY, iarray2, "my integer array");

/* gets xarray2[0] = 6.67e-11, xarray2[1] = 299792458.0 */
Util_TableGetRealArray(handle, N_XARRAY, xarray2, "my real array");
\end{verbatim}
\end{Example}
\end{ExampleSection}
\end{FunctionDescription}

%%%%%%%%%%%%%%%%%%%%%%%%%%%%%%%%%%%%%%%%

\begin{FunctionDescription}{Util\_TableSetString}
\label{Util-TableSetString}
Sets the value associated with a specified key in a table, to be
a copy of a specified C-style null-terminated character string

\begin{SynopsisSection}
\begin{Synopsis}{C}
\begin{verbatim}
#include "util_ErrorCodes.h"
#include "util_Table.h"
int status = Util_TableSetString(int handle,
                                 const char *string,
                                 const char *key);
\end{verbatim}
\end{Synopsis}
\begin{Synopsis}{Fortran}
\begin{verbatim}
call Util_TableSetString(status, handle, string, key)
integer           status, handle
character*(*)     string, key
\end{verbatim}
\end{Synopsis}
\end{SynopsisSection}

\begin{ResultSection}
\begin{ResultNote}
Results are the same as all the other \verb|Util_TableSet|* functions:
\end{ResultNote}
\begin{Result}{\rm 1}
ok (key was already in table before this call, old value was replaced)\\
   (it doesn't matter what the old value's \verb|type_code| and
    \verb|N_elements| were, \ie{} these do {\em not\/} have to match
    the new value)
\end{Result}
\begin{Result}{\rm 0}
ok (key was not in table before this call)
\end{Result}
\end{ResultSection}

\begin{ParameterSection}
\begin{Parameter}{handle ($\ge 0$)}
handle to the table
\end{Parameter}
\begin{Parameter}{string}
a pointer to the string (a C-style null-terminated string)
\end{Parameter}
\begin{Parameter}{key}
a pointer to the key (a C-style null-terminated string)
\end{Parameter}
\end{ParameterSection}

\begin{Discussion}
The key may be any C character string which does not contain a slash
character (\verb|'/'|).

The string is stored as an array of \verb|strlen(string)|
\verb|CCTK_CHAR|s.  It does {\em not\/} include a terminating
null character.

This function is very similar to \verb|Util_TableSetCharArray()|.

This function invalidates any iterators for the table which are
not in the ``null-pointer'' state.
\end{Discussion}

\begin{SeeAlsoSection}
\begin{SeeAlso}{Util\_TableCreateFromString()}
convenience routine to create a table and set key/value entries
in it based on a parameter-file--like character string
\end{SeeAlso}
\begin{SeeAlso}{Util\_TableGet*()}
get a single (1-element array) value,
or more generally the first array element of an array value
\end{SeeAlso}
\begin{SeeAlso}{Util\_TableGet*Array()}
get an array value
\end{SeeAlso}
\begin{SeeAlso}{Util\_TableGetGeneric()}
get a single (1-element array) value with generic data type
\end{SeeAlso}
\begin{SeeAlso}{Util\_TableGetGenericArray()}
get an array value with generic data type
\end{SeeAlso}
\begin{SeeAlso}{Util\_TableGet*String()}
get a character-string value
\end{SeeAlso}
\begin{SeeAlso}{Util\_TableSetCharArray()}
get an array-of-\verb|CCTK_CHAR| value
\end{SeeAlso}
\begin{SeeAlso}{Util\_TableSet*()}
set a single (1-element array) value
\end{SeeAlso}
\begin{SeeAlso}{Util\_TableSet*Array()}
set an array value
\end{SeeAlso}
\begin{SeeAlso}{Util\_TableSetGeneric()}
set a single (1-element array) value with generic data type
\end{SeeAlso}
\begin{SeeAlso}{Util\_TableSetGenericArray()}
set an array value with generic data type
\end{SeeAlso}
\begin{SeeAlso}{Util\_TableSetCharArray()}
set an array-of-\verb|CCTK_CHAR| value
\end{SeeAlso}
\end{SeeAlsoSection}

\begin{ErrorSection}
\begin{Error}{UTIL\_ERROR\_BAD\_HANDLE}
handle is invalid
\end{Error}
\begin{Error}{UTIL\_ERROR\_TABLE\_BAD\_KEY}
key contains '/' character
\end{Error}
\begin{Error}{UTIL\_ERROR\_NO\_MEMORY}
unable to allocate memory
\end{Error}
\end{ErrorSection}

\begin{ExampleSection}
\begin{Example}{C}
\begin{verbatim}
#include "util_ErrorCodes.h"
#include "util_Table.h"

static const CCTK_CHAR array[]
        = {'r', 'e', 'l', 'a', 't', 'i', 'v', 'i', 't', 'y'};
#define N_ARRAY (sizeof(array) / sizeof(array[0]))
int handle = Util_TableCreate(UTIL_TABLE_FLAGS_DEFAULT);

Util_TableSetString(handle, "relativity", "Einstein");

/* this produces the same table entry as the Util_TableSetString() */
Util_TableSetCharArray(handle, N_ARRAY, array, "Einstein");
\end{verbatim}
\end{Example}
\end{ExampleSection}
\end{FunctionDescription}

%%%%%%%%%%%%%%%%%%%%%%%%%%%%%%%%%%%%%%%%%%%%%%%%%%%%%%%%%%%%%%%%%%%%%%%%%%%%%%%%
%%%%%%%%%%%%%%%%%%%%%%%%%%%%%%%%%%%%%%%%%%%%%%%%%%%%%%%%%%%%%%%%%%%%%%%%%%%%%%%%
%%%%%%%%%%%%%%%%%%%%%%%%%%%%%%%%%%%%%%%%%%%%%%%%%%%%%%%%%%%%%%%%%%%%%%%%%%%%%%%%

\end{cactuspart}
